% !TeX root = ../../Book3.tex
% 纲要标题：## 第29章*　几何收敛、塌缩与奇异结构
% 纲要位置：几何学计划纲要.md，第 3315 行。
% 正文状态：初稿；尚未正式审核。

% 纲要内容开始
% 
% > 本章目标与证明范围：从坐标控制建立光滑预紧性，再比较 GH 和度量测度极限。曲率导数、注入半径及体积假设各有用途；所选更深正则性结论精确外引。
% 
% 纲要内容结束

\OptionalChapter{几何收敛、塌缩与奇异结构}{g3:ch:29}
\BookChapterNumber{g3:ch:29}

\begin{chapterabstract}
几何极限需要同时说明比较的对象、控制的尺度以及被保留的结构。本章由径向平行标架建立法坐标的一致导数估计，再用有限图册、固定 Euclid 嵌入和管状投影证明闭流形的光滑 Cheeger--Gromov 预紧性。圆丛与 Nil 模型展示有界曲率下的塌缩，并区分体积、直径、注入半径和维数。随后在 Ricci 下界下构造归一化度量测度极限，证明基本测度与切极限性质，精确引用非塌缩极限的结构定理。末节从松弛 Dirichlet 能量和二次输运出发，引入熵凸性与 RCD 条件，说明它们和光滑加权曲率及极限稳定性的关系。
\end{chapterabstract}

\begin{learninggoals}
	\begin{itemize}
		\item 从曲率协变导数与注入半径界取得统一法坐标，说明各项假设在光滑紧性证明中的作用。
		\item 构造局部图的共同极限与全局比较微分同胚，区分张量收敛与距离收敛。
		\item 计算圆丛塌缩的距离、体积和曲率，分析尺度变化怎样影响极限维数。
		\item 构造带测度 GH 子列，判断归一化对极限测度的影响。
		\item 准确定义正则点、切锥与奇异集，并按假设使用非塌缩 Ricci 极限的结构结论。
		\item 从松弛能量和输运熵理解合成曲率，核对光滑模型与稳定性定理的适用范围。
	\end{itemize}
\end{learninggoals}

把一个平坦环面的一条圆周不断缩短，从固定距离观察，它越来越接近另一条圆周。曲率始终为零，却有一个维度在极限中消失。若每次同时放大到短圆周的尺度，这个方向又会重新出现。另一种变化甚至不改变内在几何：把一个固定度量沿越来越复杂的微分同胚拉回，它的坐标分量可以没有有界子列。这两种情形说明，研究度量序列时既要选择合适的比较映射，也要说明在什么尺度上观察。

\ProseParagraphBreak
第10章已经用有限网和距离矩阵构造 GH 极限，但这种极限只携带距离。本章首先加入足以恢复光滑坐标的条件。曲率导数界控制坐标中的度量变化，注入半径下界保证这些坐标在统一大小的球中单射，直径上界则使有限张图足以覆盖整个流形。证明的最后一步会把极限流形与原流形真正地用微分同胚联系起来，从而能够比较联络、曲率和积分。

\ProseParagraphBreak
当统一坐标半径消失时，距离极限仍可存在。此时体积测度保留了单看距离看不到的信息：若纤维长短随位置变化，塌缩后的基底会带有非恒定的权。Ricci 下界所提供的体积比较适合控制这些极限；更深的结构定理则回答何处具有 Euclid 切锥、奇异集能有多大。最后将曲率条件写成能量与概率输运的性质，使没有光滑坐标的极限也可以进入同一类分析问题。

\ProseParagraphBreak
本章使用第2章的法坐标、第8章的体积比较、第10章的 GH 距离与预紧性，以及第11章和第21章的 Sobolev 与测度基础。光滑紧性、圆丛估计、归一化测度极限的基本构造和光滑能量比较在正文证明；一般塌缩纤维化、非塌缩 Ricci 极限结构以及合成曲率的比较与稳定性，在各自陈述处注明正式文献及所用版本。第27章的奇点放大和第28章的弱演化为这些收敛概念提供了几何背景。

% !TeX root = ../../Book3.tex
% 纲要标题：### 29.1 光滑紧性的方法
% 纲要位置：几何学计划纲要.md，第 3319 行。
% 正文状态：初稿；尚未正式审核。

% 纲要内容开始
% 
% * 谐和坐标或法坐标的统一控制
% * 过渡映射与对角子列
% * 曲率导数及注入半径条件下的紧性
% * Cheeger–Gromov 收敛与微分同胚选择
% 
% 纲要内容结束

\section{光滑紧性的方法}
\label{g3:sec:2901}

一列度量即使都平坦，其在预先指定坐标中的分量也可能剧烈振荡。寻找光滑极限时，必须连同坐标一起选择。曲率导数控制局部系数，注入半径保证这些系数定义在同样大的坐标球上，直径界则使有限张这样的图能够覆盖整个流形。本节把这三种控制逐步转成一个固定流形上的光滑收敛；特别说明局部极限怎样组成 Hausdorff 流形，以及比较用的全局微分同胚怎样构造。

\subsection{谐和坐标或法坐标的统一控制}
\label{g3:subsec:2901:01}
本节选用法坐标。\subsecref{g3:subsec:0201:04}已经证明，在法坐标中心有 \(g_{ij}=\delta_{ij}\)、\(\partial_kg_{ij}=0\)。紧性问题需要更强的量化结论：在同样大的整张坐标球上，所有度量导数都要有不依赖流形和中心的界。径向平行标架把这一问题化为沿径向线的积分方程。

\begin{lemma}\label{g3:lem:uniform-normal-coordinate-bounds}
	给定 \(n\ge2\)、\(i_0>0\) 及非负常数 \(C_0,C_1,\ldots\)。存在只依赖 \(n,i_0,C_0\) 的 \(\rho>0\)，具有以下性质。若 \((M^n,g)\) 为完备光滑黎曼流形，
	\[
		\operatorname{inj}_g(M)\ge i_0,\qquad
		|\nabla^k\mathrm{Rm}|_g\le C_k\quad(k\ge0),
	\]
	则任意中心及任意正交初始标架给出的法坐标
	\(\varphi:B(0,\rho)\subset\mathbb R^n\to M\) 满足
	\begin{equation}\label{g3:eq:uniform-normal-coordinate-bounds}
		\tfrac12 I\le (\varphi^*g)_{ij}\le2I,\qquad
		\sup_{B(0,\rho)}|\partial^\alpha(\varphi^*g)_{ij}|
		\le A_{|\alpha|}.
	\end{equation}
	每个 \(A_k\) 只依赖 \(n,i_0,C_0,\ldots,C_k\)。逆度量与联络的任意固定阶坐标导数也有相应的一致界。
\end{lemma}
\begin{proof}
	先取 \(\rho<i_0/2\)，再按曲率界缩小它。指数映射在更大的切球内可逆；\cref{g3:prop:compact-injectivity-lower-bound}中的距离结论说明 \(\varphi(B(0,\rho))=B_g(p,\rho)\)。把中心的正交基沿径向测地线平行移动，得到光滑正交标架 \(e_a(x)\)。写
	\[
		\partial_i=E_i^ae_a,\qquad
		\nabla_{\partial_i}e_b=(A_i)^a{}_b e_a,
	\]
	其中 \(E_i\) 是列向量，\(A_i\) 是反对称矩阵。光滑性来自测地线和平行移动方程对初始向量的光滑依赖；在中心 \(E_i(0)\) 为第 \(i\) 个标准基向量。径向平行性给出
	\begin{equation}\label{g3:eq:normal-radial-gauge}
		x^iA_i(x)=0,\qquad x^iE_i^a(x)=x^a.
	\end{equation}
	第二式表示径向速度在这个平行标架中的分量就是初始向量。

	设 \(\Omega_{ji}=\partial_jA_i-\partial_iA_j+[A_j,A_i]\) 是联络曲率矩阵。用第一条径向恒等式，可得
	\(x^j\Omega_{ji}=x^j\partial_jA_i+A_i\)。无挠性在此标架中为
	\(\partial_jE_i-\partial_iE_j+A_jE_i-A_iE_j=0\)；乘 \(x^j\) 并使用两条径向恒等式，得到
	\(x^j\partial_jE_i+E_i=e_i^{\mathrm{std}}+A_i x\)。沿 \(t\mapsto tx\) 积分这两个方程，便有
	\begin{equation}\label{g3:eq:normal-radial-integral-system}
		\begin{split}
			A_i(x)&=\int_0^1t x^j\Omega_{ji}(tx)\,\mathrm dt,\\
			E_i(x)&=e_i^{\mathrm{std}}+
			\int_0^1t A_i(tx)x\,\mathrm dt.
		\end{split}
	\end{equation}
	所有表达在 \(x=0\) 光滑，不需要在球心除以径向距离。

	用 \(R^a{}_{bcd}\) 表示径向正交标架中的曲率分量，则
	\(\Omega_{ji}{}^a{}_b=R^a{}_{bcd}E_j^cE_i^d\)。由
	\(x^jE_j^c(tx)=x^c\)，第一积分式中的曲率项简化为
	\[
		(A_i)^a{}_b(x)=\int_0^1t x^c
		R^a{}_{bcd}(tx)E_i^d(tx)\,\mathrm dt.
	\]
	因而在任意半径 \(r\le\rho\) 的闭子球上，记 \(\|\cdot\|_r\) 为分量上确界，有
	\[
		\|A\|_r\le c_nC_0r\|E\|_r,\qquad
		\|E-I\|_r\le c_nr\|A\|_r.
	\]
	选 \(\rho\) 使 \(c_nC_0\rho^2\) 足够小，就可吸收右端的 \(\|E\|_r\)，并使 \(\|E-I\|\) 的算子范数小于 \(1/4\)。因为 \(g_{ij}=\langle E_i,E_j\rangle\)，这给出式\eqref{g3:eq:uniform-normal-coordinate-bounds}的正定上下界。

	为取得高阶界，令 \(T_\ell\) 表示 \(\nabla^\ell\mathrm{Rm}\) 在同一正交标架中的全部分量。协变导数与分量导数满足
	\[
		\partial_iT_\ell=E_i^a(T_{\ell+1})_{a\cdots}+A_i*T_\ell,
	\]
	其中星号表示对各个标架指标的有限次线性作用，符号由张量类型确定。因此，若 \(E,A\) 到 \(k-1\) 阶的坐标导数已受控，则反复使用此式，把 \(T_0\) 的 \(k\) 阶坐标导数控制于 \(C_0,\ldots,C_k\) 及这些低阶界。

	现在对两个积分式求 \(k\) 次坐标导数。记
	\(M_k=\|\partial^kE\|_r\)、\(N_k=\|\partial^kA\|_r\)。链式法则中的 \(t\) 因子不超过一；除全部导数落在 \(E\) 或 \(A\) 上的项外，其余项均由刚才的归纳控制。所以
	\[
		N_k\le B_k+c_nC_0r M_k,\qquad
		M_k\le B'_k+c_nrN_k.
	\]
	最高阶项的系数与 \(k\) 无关，先前选择的同一个 \(\rho\) 仍能吸收。由零阶开始归纳，得到任意阶的一致界，再令 \(r\uparrow\rho\)。最后由 \(g=E^{\mathsf T}E\) 得度量导数界；微分 \(g^{-1}g=I\) 和 Christoffel 公式，得到逆度量及联络的导数界。
\end{proof}

注入半径在这里保证指数坐标没有重叠的点；曲率导数控制的则是这张坐标图中的变化率。这两个条件承担不同的工作。只知道某一点的法坐标展开，不能得到统一坐标半径，更不能控制整列流形的坐标系数。

\subsection{过渡映射与对角子列}
\label{g3:subsec:2901:02}
同一个点在相邻图中有不同坐标。若只让每张图的度量分别收敛，而没有控制图与图之间的转换，就尚未得到一个几何极限。正定上下界先控制转换的一阶导数，联络变换公式随后控制全部高阶导数。

\begin{lemma}\label{g3:lem:uniform-chart-transition-bounds}
	在上一引理的坐标图之间，过渡映射 \(y=f(x)\) 及其逆的每阶导数，在其定义域上都有由相同数据确定的上界。
\end{lemma}
\begin{proof}
	两张图的度量矩阵记为 \(G_x,G_y\)。坐标变换满足
	\[
		G_x=(\mathrm Df)^{\mathsf T}G_y(f)\mathrm Df.
	\]
	两侧的正定上下界给出 \(\mathrm Df\) 及其逆的统一范数界。联络在坐标变换下满足
	\begin{equation}\label{g3:eq:transition-second-derivative}
		\partial_i\partial_j f^a
		=\Gamma(x)^k_{ij}\partial_k f^a
		-\Gamma(y)^a_{bc}(f)\partial_i f^b\partial_j f^c.
	\end{equation}
	所以二阶导数有界。对该式反复求导，第 \(r+1\) 阶导数只用联络至 \(r-1\) 阶的导数及 \(f\) 已控制到 \(r\) 阶的导数，故逐阶归纳可得。交换两张图便得逆转换的结论。这里的估计逐点成立，不要求图的重叠域是凸集。
\end{proof}

下面把有限图册装入一个固定的 Euclid 空间。这样，极限点属于同一个 Hausdorff 空间，局部图的粘合可以直接检查。

\begin{lemma}\label{g3:lem:uniform-finite-atlas-embedding}
	在\cref{g3:lem:uniform-normal-coordinate-bounds}的假设下，进一步设 \(M\) 闭、连通且 \(\operatorname{diam}_gM\le D\)。可取固定 \(r>0\)、\(4r<\rho\)，以及至多 \(N_0\) 个法坐标图
	\[
		\varphi_i:B(0,4r)\longrightarrow M,
		\qquad M=\bigcup_{i=1}^N\varphi_i(\overline{B(0,r)}),
	\]
	其中 \(N_0\) 只依赖 \(n,D,i_0,C_0\)。由这些图可构造光滑嵌入
	\(\mathcal E:M\to\mathbb R^{N(n+1)}\)，使每个
	\(\mathcal E\circ\varphi_i\) 在 \(B(0,4r)\) 上的全部固定阶导数有一致界；其第 \(i\) 个 \((n+1)\) 维坐标块在 \(B(0,2r)\) 上恰为 \((1,x)\)。
\end{lemma}
\begin{proof}
	取一个极大的 \(r\) 分离点集 \(p_1,\ldots,p_N\)：不同点的距离至少为 \(r\)，并且不能再加点。紧性使它有限，极大性使闭 \(r\) 球覆盖 \(M\)，而开 \(r/2\) 球两两不交。法坐标中的度量下界给出每个小球的体积至少为
	\(2^{-n/2}\omega_n(r/2)^n\)。另一方面，\(|\mathrm{Rm}|\le C_0\) 蕴含 \(\operatorname{Ric}\ge-(n-1)C_0g\)，由\cref{g3:thm:bishop-gromov}及 \(M\subset B_g(p_1,D+r)\)，
	\[
		\operatorname{Vol}_gM\le V_{-C_0,n}(D+r).
	\]
	比较互不相交小球的体积和，得到所需的 \(N_0\)。

	在每个 \(p_i\) 选择正交基并取法图 \(\varphi_i\)。固定光滑径向函数 \(\chi:\mathbb R^n\to[0,1]\)，使它在 \(|x|\le2r\) 时等于一，在 \(|x|\ge3r\) 时等于零。令 \(x_i=\varphi_i^{-1}\)，把
	\[
		a_i(p)=\chi(x_i(p)),\qquad b_i(p)=a_i(p)x_i(p)
	\]
	在图外延拓为零。截断支集与 \(4r\) 图边界之间有余量，所以延拓光滑。定义
	\begin{equation}\label{g3:eq:finite-atlas-embedding}
		\mathcal E(p)=\bigl(a_1(p),b_1(p),\ldots,a_N(p),b_N(p)\bigr).
	\end{equation}
	在第 \(i\) 张图的 \(2r\) 子球上，\((a_i,b_i)=(1,x_i)\)。任一点属于某张 \(r\) 子球，故这个坐标块使 \(\mathcal E\) 的微分单射。若 \(\mathcal E(p)=\mathcal E(q)\)，取这样的 \(i\)，则 \(a_i(q)=a_i(p)=1\)；两点都在第 \(i\) 张图内，且 \(b_i(q)=b_i(p)\)，所以 \(q=p\)。紧源到 Hausdorff 空间的连续单射是到像的同胚，因此 \(\mathcal E\) 为光滑嵌入。

	最后，\(\mathcal E\circ\varphi_i\) 的每个非零坐标块是固定截断函数、坐标乘积与过渡映射的复合。上一引理控制全部导数；图外的零延拓也保持这些界。坐标块的数目有统一上界，故全映射的导数界统一。
\end{proof}

对于一列满足这些界的流形，先取子列使图数 \(N\) 固定。各个局部度量矩阵和 \(\mathcal E_j\circ\varphi_{j,i}\) 都定义在同一个 \(4r\) 球上，并有所有阶导数界。Arzelà--Ascoli 定理先对函数及一阶导数取子列，再对每阶导数依次进行；沿坐标线段积分确认极限导数仍是极限函数的导数。对有限个图和可数个阶数作对角选择，便得到在 \(\overline{B(0,3r)}\) 上同时 \(C^\infty\) 收敛的一个子列。下一步要说明这些局部极限确实描述同一个闭流形。

\subsection{曲率导数及注入半径条件下的紧性}
\label{g3:subsec:2901:03}
有限图册的嵌入保留了每张图的坐标标签：只要某块的第一分量非零，就可以用后 \(n\) 个分量除以它，恢复该图中的原坐标。这既能排除极限中出现额外重叠的片，也能把最后的局部比较推广为全局微分同胚。

\begin{theorem}[光滑紧性]\label{g3:thm:smooth-cheeger-gromov-compactness}
	设 \((M_j^n,g_j)\) 是一列闭、连通光滑黎曼流形，维数 \(n\ge2\) 固定。假设存在 \(D,i_0>0\) 和有限常数 \(C_k\)、\(k=0,1,\ldots\)，使
	\begin{equation}\label{g3:eq:smooth-compactness-assumptions}
		\operatorname{diam}_{g_j}M_j\le D,\qquad
		\operatorname{inj}_{g_j}(M_j)\ge i_0,\qquad
		\sup_{M_j}|\nabla^k\mathrm{Rm}_{g_j}|_{g_j}\le C_k.
	\end{equation}
	则存在子列、一个闭连通光滑流形 \(M_\infty^n\)、光滑黎曼度量 \(g_\infty\)，以及对该子列充分大指标定义的微分同胚
	\[
		\Psi_j:M_\infty\longrightarrow M_j,
		\qquad \Psi_j^*g_j\longrightarrow g_\infty\quad\text{于 }C^\infty(M_\infty).
	\]
\end{theorem}
\begin{proof}
	使用上一引理的共同 \(r,N_0\)，取子列使图数固定为 \(N\)，并在每张图选择正交初始标架。用同一个截断构造
	\(\mathcal E_j:M_j\to\mathbb R^{N(n+1)}\)。记
	\[
		f_{j,i}=\mathcal E_j\circ\varphi_{j,i},\qquad
		h_{j,i}=\varphi_{j,i}^*g_j.
	\]
	前一小节的对角选择给出
	\begin{equation}\label{g3:eq:compactness-local-chart-limits}
		f_{j,i}\longrightarrow f_i,\qquad
		h_{j,i}\longrightarrow h_i
		\quad\text{于 }C^\infty(\overline{B(0,3r)}),
	\end{equation}
	其中 \(\tfrac12 I\le h_i\le2I\)。以下继续用 \(j\) 表示这个子列。

	先构造极限流形。令
	\[
		S_j=\mathcal E_j(M_j),\qquad
		S=\bigcup_{i=1}^N f_i(\overline{B(0,r)})
		\subset\mathbb R^{N(n+1)}.
	\]
	因为每个 \(M_j\) 被 \(r\) 子球覆盖，
	\(S_j=\bigcup_i f_{j,i}(\overline{B(0,r)})\)。一致收敛说明 \(S_j\) 在 Euclid Hausdorff 距离下趋于非空紧集 \(S\)。对于 \(|x|<2r\)，\(f_{j,i}(x)\in S_j\)，故其极限 \(f_i(x)\) 也在 \(S\) 中；而 \(f_i\) 的第 \(i\) 块仍为 \((1,x)\)，所以它是 \(B(0,2r)\) 上的光滑嵌入。

	还须证明 \(S\) 在这张片附近没有其他分支。用 \(\alpha_i,\beta_i\) 表示环境空间第 \(i\) 块的标量和向量投影。取 \(y=f_i(x_0)\)，\(|x_0|<2r\)，并选小 \(\eta>0\)，使 \(B(x_0,2\eta)\subset B(0,2r)\)。考虑环境开集
	\[
		U=\left\{z:\alpha_i(z)>\tfrac12,
		\left|\frac{\beta_i(z)}{\alpha_i(z)}-x_0\right|<\eta\right\}.
	\]
	若 \(z\in S\cap U\)，由 Hausdorff 收敛取 \(z_j\in S_j\)、\(z_j\to z\)。充分大时 \(\alpha_i(z_j)>0\)，故 \(z_j=\mathcal E_j(p_j)\) 的原点 \(p_j\) 属于第 \(i\) 张图，且其坐标恰为
	\[
		x_j=\frac{\beta_i(z_j)}{\alpha_i(z_j)}.
	\]
	这些坐标趋于 \(x=\beta_i(z)/\alpha_i(z)\in B(x_0,\eta)\)，因而由局部收敛有 \(z=\lim f_{j,i}(x_j)=f_i(x)\)。反过来，每个 \(x\in B(x_0,\eta)\) 都给 \(f_i(x)\in S\cap U\)。所以
	\begin{equation}\label{g3:eq:limit-atlas-no-extra-sheets}
		S\cap U=f_i(B(x_0,\eta)).
	\end{equation}
	这些参数图覆盖 \(S\)，其逆直接由 \(\beta_i\) 给出。故 \(S\) 是环境空间中的光滑嵌入 \(n\) 维子流形，没有边界；它自动 Hausdorff 且具有可数基。它也连通：若紧集 \(S\) 分成两个不交非空开闭部分，两部分之间有正距离，其互不相交的小邻域将使充分接近的 \(S_j\) 也分裂，违背 \(M_j\) 连通。

	接着把局部度量 \(h_i\) 放到 \(S\) 上。若 \(f_i(x)=f_k(y)\)，\(|x|,|y|<2r\)，在这个点附近的坐标转换为
	\[
		T_{ki}(x)=\frac{\beta_k(f_i(x))}{\alpha_k(f_i(x))}.
	\]
	分母在此处等于一，缩小邻域后保持正下界。对大 \(j\)，相同表达
	\[
		T_{j,ki}(x)=\frac{\beta_k(f_{j,i}(x))}
		{\alpha_k(f_{j,i}(x))}
		=\varphi_{j,k}^{-1}\circ\varphi_{j,i}(x)
	\]
	有定义，并在较小邻域中 \(C^\infty\) 收敛到 \(T_{ki}\)。确实，恢复的坐标仍位于 \(B(0,2r)\)，所以它来自真正的第 \(k\) 张图。原度量的相容式
	\(h_{j,i}=T_{j,ki}^*h_{j,k}\) 取极限，得到
	\(h_i=T_{ki}^*h_k\)。因此存在唯一光滑正定度量 \(g_\infty\) 于 \(S\)，使 \(f_i^*g_\infty=h_i\)。

	最后构造全局比较映射。由\cref{g3:prop:euclidean-normal-tube}，紧嵌入流形 \(S\) 有固定管状邻域及光滑投影 \(\pi\) 到 \(S\)。Hausdorff 收敛使大 \(j\) 时的 \(S_j\) 位于这个邻域内，并且
	\[
		\sup_{z\in S_j}|z-\pi(z)|\longrightarrow0.
	\]
	后式也可直接由紧性证明：若有偏离零的子列，取 \(z_j\to z\in S\)，则 \(\pi(z_j)\to\pi(z)=z\)，矛盾。

	对每个 \(i\)，在 \(B(0,3r/2)\) 上令
	\[
		\nu_{j,i}=f_i^{-1}\circ\pi\circ f_{j,i}.
	\]
	这是良定义的：\(f_i(\overline{B(0,3r/2)})\) 紧含于 \(S\) 的开图 \(f_i(B(0,2r))\)，而 \(\pi\circ f_{j,i}\to f_i\) 光滑。于是
	\(\nu_{j,i}\to\operatorname{id}\) 于 \(C^\infty\)。充分大时
	\(\|\mathrm D\nu_{j,i}-I\|<1/2\)，沿凸球内的线段积分，说明 \(\nu_{j,i}\) 单射。

	对任意 \(|x|\le r\)，在闭球 \(\overline{B(x,r/4)}\) 上考虑
	\(v\mapsto x-(\nu_{j,i}(v)-v)\)。当 \(j\) 足够大时，它把这个球映入自身，且 Lipschitz 常数小于 \(1/2\)。\Cref{g3:lem:banach-contraction}给出唯一不动点，即 \(\nu_{j,i}(v)=x\)。所以 \(f_i(x)\) 至少有一个来自 \(S_j\) 的 \(\pi\) 原像。

	若 \(z\in S_j\) 满足 \(\pi(z)=f_i(x)\)、\(|x|\le r\)，则 \(|z-f_i(x)|\to0\) 一致成立。故 \(\alpha_i(z)>1/2\)，且
	\(v=\beta_i(z)/\alpha_i(z)\in B(0,3r/2)\)。恢复坐标得 \(z=f_{j,i}(v)\)，从而 \(\nu_{j,i}(v)=x\)。刚才的单射性说明原像只有一个。每个 \(S\) 中的点都属于某个闭 \(r\) 子图，故 \(\pi_j=\pi|_{S_j}:S_j\to S\) 是双射。局部表达 \(\nu_{j,i}\) 的微分可逆，所以它和其逆都是光滑的。

	取 \(M_\infty=S\)，定义
	\[
		\Psi_j=\mathcal E_j^{-1}\circ\pi_j^{-1}:S\longrightarrow M_j.
	\]
	在每个较小参数图中，
	\[
		\Psi_j\circ f_i=\varphi_{j,i}\circ\nu_{j,i}^{-1},\qquad
		f_i^*(\Psi_j^*g_j)=(\nu_{j,i}^{-1})^*h_{j,i}.
	\]
	逆函数的各阶微分公式及 \(\nu_{j,i}\to\operatorname{id}\) 使
	\(\nu_{j,i}^{-1}\to\operatorname{id}\) 光滑。结合局部度量收敛，右端趋于 \(h_i=f_i^*g_\infty\)。有限个核心图覆盖 \(S\)，故得到全流形上的 \(C^\infty\) 收敛。
\end{proof}

证明中使用的 Euclid 嵌入只负责编码坐标，它不要求保持黎曼度量。极限度量由 \(h_i\) 的相容性另行定义。管状投影的单射性也单独作了检验；只有局部微分可逆，尚不能排除一个多层覆盖。

\begin{corollary}\label{g3:cor:bounded-geometry-diffeomorphism-finiteness}
	在固定 \(n,D,i_0,(C_k)_{k\ge0}\) 的式\eqref{g3:eq:smooth-compactness-assumptions}所描述的闭连通流形类中，只有有限多个微分同胚类型。
\end{corollary}
\begin{proof}
	若有无穷多个类型，从每种类型选一个带有所需度量的流形。紧性定理给出一个子列及到同一个 \(M_\infty\) 的微分同胚，因此子列中充分靠后的流形类型相同，与选择矛盾。
\end{proof}

\subsection{Cheeger–Gromov 收敛与微分同胚选择}
\label{g3:subsec:2901:04}
上述极限并不把原来各流形上的点预先认作同一点。真正用来比较的是微分同胚的拉回；这一点也是光滑几何收敛的定义所要记录的内容。

\begin{definition}\label{g3:def:cheeger-gromov-convergence}
	设 \((M_j,g_j)\)、\((M_\infty,g_\infty)\) 为闭光滑黎曼流形，维数相同，\(k\ge0\) 为整数。若存在对充分大 \(j\) 定义的微分同胚
	\(\Psi_j:M_\infty\to M_j\)，使
	\[
		\Psi_j^*g_j\longrightarrow g_\infty\quad\text{于 }C^k(M_\infty),
	\]
	则称这列流形 \(C^k\) \term[Cheeger Gromov shoulian]{Cheeger--Gromov 收敛}[Cheeger--Gromov convergence]到 \((M_\infty,g_\infty)\)。若同一组 \(\Psi_j\) 对每个有限阶 \(k\) 都给出收敛，则称其光滑收敛或 \(C^\infty\) 收敛。范数可用极限流形上的任意固定光滑度量及有限图册定义，紧性使不同选择给出相同的收敛含义。
\end{definition}

\begin{example}\label{g3:exm:coordinate-oscillation-smooth-limit}
	在平坦二维环面 \(\mathbb T^2=(\mathbb R/2\pi\mathbb Z)^2\) 上取 \(g_0=\mathrm dx^2+\mathrm dy^2\)，令
	\(\phi_j(x,y)=(x+j\sin y,y)\)，并置
	\[
		g_j=\phi_j^*g_0=(\mathrm dx+j\cos y\,\mathrm dy)^2+\mathrm dy^2.
	\]
	每个 \(\phi_j\) 都是微分同胚。各 \(g_j\) 的曲率及其导数全为零，直径与注入半径也相同；但固定坐标中的 \(\mathrm dy^2\) 系数在 \(y=0\) 等于 \(1+j^2\)，没有有界子列。选择 \(\Psi_j=\phi_j^{-1}\) 后，\(\Psi_j^*g_j=g_0\) 恒成立，所以它们光滑 Cheeger--Gromov 收敛。这里变化的是标记，内在几何从未变化。
\end{example}

\begin{proposition}\label{g3:prop:smooth-convergence-implies-gh}
	若一列闭连通光滑黎曼流形 \((M_j,g_j)\) 以 \(C^0\) Cheeger--Gromov 意义收敛到闭连通 \((M,g)\)，则其黎曼距离空间以 Gromov--Hausdorff 意义收敛到 \((M,d_g)\)。此外总体积收敛；对任意连续函数 \(u\) 于 \(M\)，
	\[
		\int_Mu\,\mu_{\Psi_j^*g_j}\longrightarrow\int_Mu\,\mu_g.
	\]
\end{proposition}
\begin{proof}
	记 \(h_j=\Psi_j^*g_j\)。紧性、\(g\) 正定及 \(C^0\) 收敛给出 \(\varepsilon_j\downarrow0\)，使
	\((1-\varepsilon_j)g\le h_j\le(1+\varepsilon_j)g\)。对所有分段光滑曲线积分长度，再取路径下确界，得到
	\[
		\sqrt{1-\varepsilon_j}\,d_g(x,y)
		\le d_{h_j}(x,y)
		\le\sqrt{1+\varepsilon_j}\,d_g(x,y).
	\]
	因 \(\operatorname{diam}_gM<\infty\)，两距离之差一致趋于零。\(\Psi_j\) 的图像是 \(M\) 与 \(M_j\) 的一个满对应，其畸变正是这些距离差的上确界；\cref{g3:thm:gh-correspondence}遂给出 GH 收敛。

	在固定坐标中，\(\mu_{h_j}/\mu_g=\sqrt{\det(g^{-1}h_j)}\to1\) 一致成立。与有界连续 \(u\) 相乘后积分，便得测度结论，取 \(u=1\) 即得总体积收敛。
\end{proof}

对于没有统一直径界的完备流形，通常只观察一个基点附近越来越大的区域。相应的\term[daijidian guanghua shoulian]{带基点光滑收敛}[pointed smooth convergence]使用完备连通的 \((M_j,g_j,p_j)\) 与 \((M_\infty,g_\infty,p_\infty)\)：取包含 \(p_\infty\) 的开集穷尽 \(U_j\)，以及光滑嵌入
\[
	\Psi_j:U_j\longrightarrow M_j,\qquad
	\Psi_j(p_\infty)=p_j,
\]
使 \(M_\infty\) 的每个紧集上的拉回度量都光滑收敛，并使像包含半径 \(R_j\to\infty\) 的基点球 \(B_{g_j}(p_j,R_j)\)。\(C^k\) 版本把这里的光滑收敛改为紧集上的 \(C^k\) 收敛。这个定义同时指定了比较区域和覆盖范围；它允许极限非紧，而本节的紧性定理给出的是有统一直径界时的闭流形版本。

\ProseParagraphBreak
光滑收敛携带度量张量、联络及曲率的局部信息，GH 收敛只比较距离。上面的命题给出了由前者到后者的方向。若注入半径趋于零，即使曲率始终有界，局部坐标也可能无法保持同样的有效半径；下一节的平坦环面和圆丛模型将展示这种维数变化。
% 纲要条目排为编号小节标题，并进入目录与书签；说明段仍保留于纲要注释中。

% !TeX root = ../../Book3.tex
% 纲要标题：### 29.2 塌缩与维数变化
% 纲要位置：几何学计划纲要.md，第 3326 行。
% 正文状态：初稿；尚未正式审核。

% 纲要内容开始
% 
% * 扁平环面和圆丛塌缩
% * 体积、直径和曲率的不同控制
% * 局部纤维结构的深层定理接口
% * 极限维数与拓扑维数
% 
% 纲要内容结束

\section{塌缩与维数变化}
\label{g3:sec:2902}

第10章已经算过扁平二维环面趋于圆周的距离误差。现在要观察同一过程中更多的量：体积怎样消失，注入半径为何变小，局部曲率是否仍能保持有界，以及缩短的方向在拓扑上是否总能拼成直积。先把圆周方向的缩短推广到一般圆丛，再用一个非平凡的三维模型看清纤维的扭转。

\subsection{扁平环面和圆丛塌缩}
\label{g3:subsec:2902:01}
设 \((B^m,h)\) 为闭连通黎曼流形，\(P\to B\) 为一个光滑主圆丛，圆群取 \(\mathbb R/(2\pi\mathbb Z)\)。令 \(V\) 为周期 \(2\pi\) 的基本向量场，选圆作用不变的实联络一形式 \(\theta\)，满足 \(\theta(V)=1\)。水平空间为 \(\ker\theta\)，其上投影微分与基底切空间同构。对 \(0<\varepsilon\le1\)，定义
\begin{equation}\label{g3:eq:circle-bundle-collapsing-metric}
	g_\varepsilon=\pi^*h+\varepsilon^2\theta\otimes\theta.
\end{equation}
这里第一项在水平空间上正定、在垂直空间上为零，第二项恰好补上垂直方向的长度。每条圆纤维的周长为 \(2\pi\varepsilon\)。这样的度量称为\term[lianluo duliang]{联络度量}[connection metric]。

\begin{proposition}\label{g3:prop:circle-bundle-collapse}
	在上述闭基底、光滑主圆丛与固定联络的条件下，对任意 \(p,q\in P\)，有
	\begin{equation}\label{g3:eq:circle-bundle-distance-bound}
		d_h(\pi p,\pi q)\le d_{g_\varepsilon}(p,q)
		\le d_h(\pi p,\pi q)+\pi\varepsilon.
	\end{equation}
	因此
	\[
		d_{\mathrm{GH}}((P,g_\varepsilon),(B,h))\le\frac{\pi\varepsilon}{2},
		\qquad \operatorname{Vol}_{g_\varepsilon}P
		=2\pi\varepsilon\operatorname{Vol}_hB.
	\]
	每个 \(p\in P\) 的注入半径不超过 \(\pi\varepsilon\)。此外，存在只依赖固定 \((B,h,\theta)\) 的常数 \(C\)，使 \(\sup_P|\mathrm{Rm}_{g_\varepsilon}|\le C\) 对 \(0<\varepsilon\le1\) 成立。
\end{proposition}
\begin{proof}
	投影微分在水平向量上保长，在垂直向量上为零，所以投影不增加道路长度，给出第一个距离不等式。反过来，将从 \(\pi p\) 到 \(\pi q\) 的最短线从 \(p\) 水平提升。提升是一条与原曲线等长的道路，其终点与 \(q\) 在同一圆纤维上；再沿这条纤维的一段短弧前进，长度至多 \(\pi\varepsilon\)。水平提升在整个区间存在，因为它是紧丛 \(P\) 上光滑常微分方程的解。两段相接得到第二个不等式。投影为满射，它的图是畸变至多为 \(\pi\varepsilon\) 的对应，故\cref{g3:thm:gh-correspondence}给出 GH 估计。

	在局部平凡化中写 \(\theta=\mathrm d\varphi+A\)，其中 \(\varphi\) 为圆周变量。以水平正交余标架连同 \(\varepsilon\theta\) 计算体积密度，得到 \(\mu_{g_\varepsilon}=\varepsilon\,\mathrm d\varphi\,\mu_h\)。沿圆纤维积分即得体积公式。以下联络计算还表明纤维是测地线；从 \(p\) 出发沿两个相反纤维方向各行长度 \(\pi\varepsilon\)，到达同一点。指数映射因而不可能在更大的切球上单射，给出注入半径上界。

	为核对曲率界，令 \(F\) 为基底上的二形式，使 \(\mathrm d\theta=\pi^*F\)。在一个固定基底坐标片中取正交场 \(e_1,\ldots,e_m\)，以 \(E_i\) 表示其不变水平提升，令 \(E_0=\varepsilon^{-1}V\)。则 \(E_0,E_1,\ldots,E_m\) 为 \(g_\varepsilon\) 正交标架，并且
	\[
		[E_i,E_j]=[e_i,e_j]^H-\varepsilon F_{ij}E_0,
		\qquad [E_i,E_0]=0.
	\]
	Koszul 公式逐项给出
	\begin{equation}\label{g3:eq:circle-bundle-connection}
	\begin{aligned}
		\nabla_{E_i}E_j&=(\nabla^h_{e_i}e_j)^H-\tfrac\varepsilon2F_{ij}E_0,\\
		\nabla_{E_i}E_0=\nabla_{E_0}E_i&=\tfrac\varepsilon2\sum_jF_{ij}E_j,
		&\nabla_{E_0}E_0&=0.
	\end{aligned}
	\end{equation}
	所有系数都是基底光滑函数的提升；沿 \(E_0\) 的导数为零，沿 \(E_i\) 的导数就是沿 \(e_i\) 的导数。曲率由联络系数的一阶导数、两个联络系数的乘积及括号系数组成。因此在基底的有限个相对紧标架片内，每个正交曲率分量都是有界基底系数与 \(\varepsilon\)、\(\varepsilon^2\) 的有限组合，没有负次幂。这给出与 \(\varepsilon\) 无关的统一界，并证明了前面用到的纤维测地性。
\end{proof}

这个计算解释了有界曲率与小注入半径为何能够共存。缩短圆周并没有产生大小为 \(\varepsilon^{-2}\) 的圆周内在曲率，因为一维纤维本身平坦；其与水平空间的耦合则在正交标架中带着正次幂 \(\varepsilon\)。

\begin{example}\label{g3:exm:nil-circle-collapse}
	取\cref{g3:prop:ricci-flow-nil-model}所用的群乘法
	\[
		(x,y,z)(x',y',z')=(x+x',y+y',z+z'+xy')
	\]
	及整数子群 \(\Gamma\)，在紧商 \(P=\Gamma\backslash\mathbb R^3\) 上令
	\[
		g_\varepsilon=\mathrm dx^2+\mathrm dy^2+
		\varepsilon^2(\mathrm dz-x\,\mathrm dy)^2.
	\]
	投影到 \((x\bmod1,y\bmod1)\) 给出到单位平坦二维环面的圆丛，纤维周期在 \(z\) 坐标中为一。整数左平移保持三个所列一形式，故度量和投影在商上良定义。前一命题的距离论证在此给
	\[
		d_{\mathrm{GH}}((P,g_\varepsilon),\mathbb T^2)\le\varepsilon/4,
		\qquad \operatorname{Vol}_{g_\varepsilon}P=\varepsilon.
	\]
	第27章已经从 Koszul 公式算出这一度量在水平与垂直基本平面上的截面曲率为
	\[
		K_{12}=-3\varepsilon^2/4,\qquad
		K_{13}=K_{23}=\varepsilon^2/4.
	\]
	因而所有曲率分量趋于零，而整个三维流形趋于二维环面。

	这个圆丛不是直积圆丛。若存在全局光滑截面 \(\sigma\)，由 \(\gamma=\mathrm dz-x\,\mathrm dy\) 得
	\(\mathrm d(\sigma^*\gamma)=-\mathrm dx\wedge\mathrm dy\)。对基底环面积分，左边由 Stokes 定理为零，右边为 \(-1\)，矛盾。圆纤维在每个小坐标片里都可见，但其整体拼接保留了非平凡的扭转。
\end{example}

\subsection{体积、直径和曲率的不同控制}
\label{g3:subsec:2902:02}
光滑紧性使用统一的坐标半径及系数控制。直径限制整个空间的尺度；体积衡量总体质量；曲率控制局部弯曲。三者各自有界，并不表示它们可以相互替代。

\begin{example}\label{g3:exm:collapse-control-counterexamples}
	以下各族都由闭连通光滑流形组成，距离均为内在黎曼距离。
	在 \(T_{1,\varepsilon}=S^1_1\times S^1_\varepsilon\) 上，曲率恒为零，直径为 \(\tfrac12\sqrt{1+\varepsilon^2}\)，面积为 \(\varepsilon\)，注入半径为 \(\varepsilon/2\)（\(0<\varepsilon\le1\)）。因此曲率和直径界允许面积及注入半径趋零。

	\ProseParagraphBreak
	在 \(T_{L,L^{-1}}\)、\(L\to\infty\)，上，面积恒为一、曲率仍为零，但直径至少为 \(L/2\)。故总体积下界配合曲率界也不控制整体直径。在半径 \(\varepsilon\) 的圆球上，直径为 \(\pi\varepsilon\)、体积与 \(\varepsilon^n\) 成正比，而截面曲率为 \(\varepsilon^{-2}\)。这个缩球趋于一点时，失败的是统一曲率界。

	\ProseParagraphBreak
	另取周期 \(2\pi\) 的二维平坦环面，令
	\[
		u_j(x,y)=j^{-1}\sin(jx),\qquad
		g_j=e^{2u_j}(\mathrm dx^2+\mathrm dy^2).
	\]
	有 \(e^{-2/j}g_0\le g_j\le e^{2/j}g_0\)，故恒等映射的距离畸变趋零，体积也趋于 \(\operatorname{Vol}_{g_0}\)。但是\cref{g3:cor:conformal-scalar-power}的二维公式给
	\[
		K_{g_j}=e^{-2u_j}\Delta_{g_0}u_j
		=j e^{-2u_j}\sin(jx),
	\]
	其绝对值上确界趋于无穷。距离可以收敛到同一个光滑空间，同时二阶曲率没有统一控制。
\end{example}

\begin{proposition}\label{g3:prop:collapse-local-global-volume}
	固定 \(n\ge2\)、\(k\in\mathbb R\) 及 \(D>0\)。设闭连通 \(n\) 维黎曼流形 \((M_j,g_j)\) 满足 \(\operatorname{Ric}_{g_j}\ge(n-1)kg_j\)、\(\operatorname{diam}M_j\le D\)。则
	\[
		\operatorname{Vol}M_j\to0
		\quad\Longleftrightarrow\quad
		\inf_{p\in M_j}\operatorname{Vol}B(p,1)\to0.
	\]
	更精确地，存在只依赖 \(n,k,D\) 的 \(c>0\)，使每个 \(p\in M_j\) 都满足
	\[
		c\operatorname{Vol}M_j\le\operatorname{Vol}B(p,1)
		\le\operatorname{Vol}M_j.
	\]
\end{proposition}
\begin{proof}
	置 \(k_0=\min\{k,0\}\)、\(R=2D+2\)，则 \(B(p,R)=M_j\)。\Cref{g3:thm:bishop-gromov}的体积比较给
	\[
		\frac{\operatorname{Vol}B(p,1)}{V_{k_0,n}(1)}
		\ge\frac{\operatorname{Vol}B(p,R)}{V_{k_0,n}(R)}.
	\]
	取 \(c=V_{k_0,n}(1)/V_{k_0,n}(R)>0\) 即得下界；上界来自集合包含。对各点取下确界，两个方向随即成立。
\end{proof}

在这组固定维数、Ricci 下界和直径条件下，总体积趋零因而也表现为每个固定单位尺度的体积趋零。相对地，\secref{g3:sec:2901}的法坐标控制与注入半径正下界提供一个统一小半径 \(r\)，使拉回度量至少为 Euclid 度量的一半，于是每个 \(r\)-球体积至少为 \(2^{-n/2}\omega_nr^n\)。圆丛塌缩恰好失去这个统一坐标半径。

\begin{table}[htbp]
	\centering\small
	\caption{四个度量族中不同几何控制的表现}
	\label{g3:tab:collapse-controls}
	\begin{tabularx}{\textwidth}{@{}lXXX@{}}
		\toprule
		度量族 & 曲率 & 总体积 & 直径或极限\\
		\midrule
		\(T_{1,\varepsilon}\) & 恒为零 & 趋于零 & 直径有界，极限为圆周\\
		\(T_{L,L^{-1}}\) & 恒为零 & 恒为一 & 直径趋于无穷\\
		半径 \(\varepsilon\) 的球面 & 无统一上界 & 趋于零 & 直径趋零，极限为一点\\
		\(e^{2j^{-1}\sin(jx)}g_0\) & 无统一绝对值界 & 趋于正数 & GH 极限为原平坦环面\\
		\bottomrule
	\end{tabularx}
\end{table}

\subsection{局部纤维结构的深层定理接口}
\label{g3:subsec:2902:03}
圆丛算例预先给出了投影；一般塌缩序列只给出距离近似。能否把这些近似改造成光滑纤维化，是一个额外的几何结论。即使存在纤维，纤维也可能是更高维的幂零空间，而不只是圆或环面。

\begin{definition}\label{g3:def:infranilmanifold}
	设 \(N\) 为连通单连通幂零 Lie 群，\(\operatorname{Aff}(N)=N\rtimes\operatorname{Aut}(N)\)。设离散子群 \(\Gamma\subset\operatorname{Aff}(N)\) 自由、正常不连续且余紧地作用于 \(N\)，并满足 \([\Gamma:\Gamma\cap N]<\infty\)。紧光滑商 \(\Gamma\backslash N\) 称为\term[infranil liuxing]{infranil 流形}[infranilmanifold]。若 \(\Gamma\subset N\)，则称为\term[miling liuxing]{幂零流形}[nilmanifold]。
\end{definition}

幂零是指 Lie 代数的下降中心列在有限步后变为零。有限指标条件说明，infranil 流形有一个有限覆盖是幂零流形；它不要求给定黎曼度量平坦。\(\mathbb R^r/\mathbb Z^r\) 是最基本的例子，上面算过的三维 Nil 紧商则具有非交换的群结构。Fukaya\footnote{\OriginalHanName{深谷賢治}（ふかや けんじ，Kenji Fukaya，1959--），日本数学家，研究黎曼流形的塌缩与辛几何；姓名读音见\href{https://www.japan-acad.go.jp/japanese/members/4/fukaya_kenji.html}{日本学士院会员资料}。}的纤维化定理把这些模型推广到具有统一曲率界的塌缩序列。

\begin{theorem}[Fukaya 纤维化定理]\label{g3:thm:fukaya-fibration-external}
	设 \((M_j^n,g_j)\) 是一列闭连通光滑黎曼流形，\((B^m,h)\) 为固定闭连通光滑黎曼流形，\(1\le m<n\)。若
	\[
		|\sec_{g_j}|\le1,\qquad
		d_{\mathrm{GH}}((M_j,g_j),(B,h))\longrightarrow0,
	\]
	则对充分大的 \(j\)，存在光滑局部平凡纤维丛 \(\pi_j:M_j\to B\)，其纤维为 \(n-m\) 维 infranil 流形。
\end{theorem}

本定理作为外部结果使用，取自\cite[定理0-1，第333页的(0-3-1)、(0-3-2)]{Fukaya1989CollapsingII}；同文注0-4说明这两项已在作者的前一篇工作中证明。这里选取的是光滑闭极限上的纤维化版本。它在基底的每个充分小邻域内给出光滑乘积表示，并把局部表示的整体拼接保留为丛；所列双侧截面曲率界与光滑基底条件都是该陈述的输入。Ricci 下界极限以及带奇异点的基底将在下一节用另一组结构结论处理。

\ProseParagraphBreak
当 \(n-m=1\) 时，紧连通的一维 infranil 流形是圆，因而恢复圆纤维的图景。纤维丛本身仍可能扭转，\cref{g3:exm:nil-circle-collapse}正展示了这一点。若只是总直径趋零而把空间压成一点，也不能直接把一个任意流形称为 infranil 流形：缩球算例的曲率已经发散，失去了本定理的假设。

\subsection{极限维数与拓扑维数}
\label{g3:subsec:2902:04}
“维数下降”在光滑模型中容易看见，讨论一般极限时却应先说明所用的维数。光滑流形的坐标维数由局部坐标确定，Hausdorff 维数由距离确定；拓扑维数则只使用开集结构。这些定义使我们能够分别追踪距离改变和拓扑改变。

\ProseParagraphBreak
第21章的 Hausdorff 维数由覆盖集合的直径幂确定：
\[
	\dim_H X=\inf\{s\ge0:\mathcal H^s(X)=0\}.
\]
拓扑上，对紧度量空间 \(X\)，若任意有限开覆盖都有一个有限开加细，使每点至多同时属于 \(m+1\) 个加细开集，就说其\term[fugai weishu]{覆盖维数}[covering dimension]不超过 \(m\)，记为 \(\dim_{\mathrm{top}}X\le m\)。开加细是指覆盖中的每个新开集都包含在某个原开集中。把开覆盖及其加细沿同胚搬运，覆盖重数不变，故这个维数确为同胚不变量。下面在区间上直接比较它与 Hausdorff 维数。

\begin{proposition}\label{g3:prop:gh-snowflake-dimension}
	对任意 \(0<\alpha<1\)，区间 \([0,1]\) 上的函数 \(d_\alpha(x,y)=|x-y|^\alpha\) 是距离，诱导通常拓扑，而
	\[
		\dim_H([0,1],d_\alpha)=\alpha^{-1}.
	\]
	其覆盖维数为一。另令 \(E_j=\{k/j:0\le k\le j\}\) 并赋予通常距离，则 \(E_j\) 在 GH 意义下趋于通常区间，但每个 \(E_j\) 的 Hausdorff 维数和覆盖维数都为零。
\end{proposition}
\begin{proof}
	因 \((a+b)^\alpha\le a^\alpha+b^\alpha\)，三角不等式成立；\(d_\alpha\)-球恰为通常半径取 \(1/\alpha\) 次幂的球，故拓扑不变。对任意集合 \(U\)，有 \(\operatorname{diam}_{d_\alpha}U=(\operatorname{diam}_{|\cdot|}U)^\alpha\)。因此 Hausdorff 覆盖和中每项的 \(s\) 次幂，恰为通常距离下的 \(\alpha s\) 次幂，覆盖尺度也同步变换，临界指数变为 \(1/\alpha\)。这里通常区间的临界指数为一：等分小区间给出维数至多一；任意覆盖的区间包络长度之和至少为一，给出维数至少一。

	对区间的任意有限开覆盖，紧性给出 Lebesgue 数 \(\delta>0\)：直径小于 \(\delta\) 的子集包含在一个覆盖开集中。取充分细的等距网点，以间距的 \(3/5\) 为半径取相对开区间。这些小区间覆盖 \([0,1]\)，每点至多属于两个，且直径小于 \(\delta\)，所以给出重数至多二的开加细。故覆盖维数至多一。若它为零，则覆盖 \([0,2/3)\)、\((1/3,1]\) 有互不相交的有限开加细；这些开集不能只有一个，却将连通区间分成不交的非空开集，矛盾。故维数等于一；雪花距离的拓扑相同，覆盖维数也相同。

	有限集在任何正次幂的 Hausdorff 测度下为零，可用单点覆盖；其离散拓扑也使任何开覆盖都能加细为重数一的单点覆盖。因此两种维数均为零。另一方面，\(E_j\subset[0,1]\) 且每个区间点距它至多 \(1/(2j)\)，所以共同环境中的 Hausdorff 距离趋零，进而 GH 收敛。
\end{proof}

雪花距离改变了 Hausdorff 维数而保留拓扑；有限网逼近则说明，任意 GH 序列的 Hausdorff 维数也可以在极限中增加。对于统一维数、Ricci 下界的黎曼序列，必须进一步利用曲率和测度信息，才能得到更强的维数结论。

\ProseParagraphBreak
对圆丛塌缩，每个 \(P\) 都是 \((m+1)\) 维光滑流形，极限 \(B\) 是 \(m\) 维；这是在原尺度观察时的维数变化。若放大到纤维尺度，缩短的方向可以重新出现。\Cref{g3:prop:torus-collapse-rescale}已准确计算：当扁平环面的纤维周长为 \(h_j\)、观察尺度为 \(r_j\) 时，\(h_j/r_j\) 趋零、趋于正数或趋于无穷，分别产生直线、圆柱或平面。极限对象既取决于原度量，也取决于比较空间时选定的尺度。
% 纲要条目排为编号小节标题，并进入目录与书签；说明段仍保留于纲要注释中。

% !TeX root = ../../Book3.tex
% 纲要标题：### 29.3 Ricci 下界的极限
% 纲要位置：几何学计划纲要.md，第 3333 行。
% 正文状态：初稿；尚未正式审核。

% 纲要内容开始
% 
% * 体积比较及非塌缩条件
% * 测度的归一化和收敛
% * 正则点、切锥及奇异集
% * 所选结构定理的准确范围
% 
% 纲要内容结束

\section{Ricci 下界的极限}
\label{g3:sec:2903}

Ricci 下界控制同心球的体积比，因而即使总体积趋于零，也能控制一只球内可以放下多少个互不相交的小球。这给出距离空间的紧性。为了继续讨论积分，还要保留体积在极限中的分布：将总体积归一化后，缩短的纤维仍会通过权函数影响极限测度。本节先证明这两层紧性，再说明非塌缩假设怎样使极限测度恢复为相应维数的 Hausdorff 测度。

\subsection{体积比较及非塌缩条件}
\label{g3:subsec:2903:01}
固定整数 \(n\ge2\)、常数 \(K\ge0\) 和 \(D>0\)，考虑连通闭黎曼流形 \((M^n,g)\)，满足
\begin{equation}\label{g3:eq:ricci-limit-class}
	\operatorname{Ric}_g\ge-(n-1)Kg,
	\qquad \operatorname{diam}(M,g)\le D.
\end{equation}
这里“闭”仍表示紧且无边界。沿用\secref{g3:sec:0802}的记号，\(V_{-K,n}(r)\) 是常截面曲率 \(-K\) 的单连通模型中半径 \(r\) 的球体积；特别地，\(V_{0,n}(r)=\omega_nr^n\)，其中 \(\omega_n\) 为 Euclid 单位球体积。

\begin{proposition}\label{g3:prop:ricci-limit-volume-packing}
	设连通闭黎曼流形 \((M^n,g)\) 满足式\eqref{g3:eq:ricci-limit-class}，并令 \(R_0=2D+1\)。则对任意 \(x\in M\) 和 \(0<r<R_0\)，有
	\begin{equation}\label{g3:eq:ricci-limit-relative-volume}
		\frac{\operatorname{Vol}_gB(x,r)}{\operatorname{Vol}_gM}
		\ge\frac{V_{-K,n}(r)}{V_{-K,n}(R_0)},
		\qquad \operatorname{Vol}_gM\le V_{-K,n}(R_0).
	\end{equation}
	对 \(0<\varepsilon<R_0\)，以半径 \(\varepsilon\) 的开球覆盖 \(M\) 所需的最少球数满足
	\begin{equation}\label{g3:eq:ricci-limit-packing}
		N_M(\varepsilon)
		\le\frac{V_{-K,n}(R_0)}{V_{-K,n}(\varepsilon/2)}.
	\end{equation}
	若另外有 \(\operatorname{Vol}_gM\ge v>0\)，则
	\begin{equation}\label{g3:eq:ricci-limit-noncollapse-balls}
		\operatorname{Vol}_gB(x,r)
		\ge v\frac{V_{-K,n}(r)}{V_{-K,n}(R_0)}.
	\end{equation}
\end{proposition}
\begin{proof}
	因为 \(R_0>D\)，故 \(B(x,R_0)=M\)。\cref{g3:thm:bishop-gromov}给出
	\[
		\frac{\operatorname{Vol}_gB(x,r)}{V_{-K,n}(r)}
		\ge\frac{\operatorname{Vol}_gM}{V_{-K,n}(R_0)},
	\]
	以及模型球体积上界，得到式\eqref{g3:eq:ricci-limit-relative-volume}。

	取一个极大有限 \(\varepsilon\)-分离点集 \(\{x_1,\ldots,x_N\}\)，其中不同点距离至少为 \(\varepsilon\)。紧性保证这种点集有限，极大性保证半径 \(\varepsilon\) 的开球覆盖 \(M\)。半径 \(\varepsilon/2\) 的开球两两不交，所以
	\[
		N\frac{V_{-K,n}(\varepsilon/2)}{V_{-K,n}(R_0)}
		\operatorname{Vol}_gM
		\le\sum_{a=1}^N\operatorname{Vol}_gB(x_a,\varepsilon/2)
		\le\operatorname{Vol}_gM.
	\]
	消去正的总体积即得覆盖数界。最后把总体积下界代入相对体积估计。
\end{proof}

这里恢复的是\cref{g3:cor:ricci-diameter-gh-precompact}中的关键机制：小球数估计所用的总体积最终被消去了。因此 GH 预紧性无需假设体积下界；非塌缩的作用是为极限保留原来的 \(n\) 维体积。

\begin{definition}\label{g3:def:ricci-limit-noncollapse}
	对满足式\eqref{g3:eq:ricci-limit-class}的连通闭 \(n\) 维黎曼流形序列，若存在 \(v>0\) 使所有 \(i\) 都有 \(\operatorname{Vol}_{g_i}M_i\ge v\)，就称该序列在此尺度下\term[feitansu ricci xulie]{非塌缩}[noncollapsed]。若 \(\operatorname{Vol}_{g_i}M_i\to0\)，就称它在此尺度下体积塌缩。
\end{definition}

在统一直径界下，这个非塌缩条件等价于存在基点 \(p_i\) 和 \(v_1>0\)，使 \(\operatorname{Vol}_{g_i}B(p_i,1)\ge v_1\)。一个方向由球包含于全空间得到，另一方向由式\eqref{g3:eq:ricci-limit-noncollapse-balls}在 \(r=1\) 处得到。对一般完备非紧序列，总体积可能无穷，此时通常直接以基点单位球的体积下界定义非塌缩。

\subsection{测度的归一化和收敛}
\label{g3:subsec:2903:02}
紧致情形最方便的归一化是
\begin{equation}\label{g3:eq:ricci-probability-normalization}
	\nu_i=\frac{\mathrm d\operatorname{vol}_{g_i}}
	{\operatorname{Vol}_{g_i}M_i}.
\end{equation}
它把每个 \(M_i\) 看成总质量为一的空间。要比较不同空间上的积分，先把它们放进同一个紧致环境，再使用固定空间上的测度弱收敛。

\begin{lemma}\label{g3:lem:compact-probability-subsequence}
	设 \(Z\) 为非空紧度量空间，\(\mu_i\) 为 \(Z\) 上的 Borel 概率测度。则存在子列和 Borel 概率测度 \(\mu\)，使
	\begin{equation}\label{g3:eq:compact-measure-weak-convergence}
		\int_Z f\,\mathrm d\mu_{i_j}\longrightarrow
		\int_Z f\,\mathrm d\mu\qquad(f\in C(Z)).
	\end{equation}
	这称为概率测度的弱收敛，记作 \(\mu_{i_j}\rightharpoonup\mu\)。对于闭集 \(F\subset Z\) 和开集 \(G\subset Z\)，分别有
	\begin{equation}\label{g3:eq:weak-measure-closed-open}
		\limsup_j\mu_{i_j}(F)\le\mu(F),
		\qquad \mu(G)\le\liminf_j\mu_{i_j}(G).
	\end{equation}
\end{lemma}
\begin{proof}
	先构造 \(C(Z)\) 中的可数一致稠密集。对每个正整数 \(m\)，取有限 \(1/m\)-网 \(z_{m,a}\)，并令
	\[
		\psi_{m,a}(z)=
		\frac{(2/m-d(z,z_{m,a}))_+}
		{\sum_b(2/m-d(z,z_{m,b}))_+}.
	\]
	分母处处为正，这些连续函数非负且和为一。所有函数 \(\sum_aq_a\psi_{m,a}\)，其中 \(m\) 为正整数、\(q_a\) 为有理数，组成可数集。给定连续函数 \(f\)，取 \(q_a\) 逼近 \(f(z_{m,a})\)，其一致误差由 \(f\) 在距离 \(2/m\) 内的振幅和有理数逼近误差控制，故此集一致稠密。

	将这个稠密集列成 \(f_1,f_2,\ldots\)。由于 \(|\int f_a\,\mathrm d\mu_i|\le\|f_a\|_\infty\)，逐个选子列再取对角子列，使这些积分全部收敛。一致逼近随后表明，对每个 \(f\in C(Z)\)，积分也收敛，其极限 \(L(f)\) 是正的有界线性泛函，且 \(L(1)=1\)。\subsecref{g3:subsec:2101:02}说明的连续函数空间测度表示定理给出唯一正 Radon 测度 \(\mu\)，满足 \(L(f)=\int f\,\mathrm d\mu\)；\(L(1)=1\) 保证它是概率测度。

	对非空闭集 \(F\)，连续函数
	\(f_m(z)=(1-md(z,F))_+\) 随 \(m\to\infty\) 递减到 \(\mathbf1_F\)。因此
	\[
		\limsup_j\mu_{i_j}(F)
		\le\int_Zf_m\,\mathrm d\mu
		\longrightarrow\mu(F).
	\]
	空集情形直接成立。对补集使用概率测度的总质量为一，得到开集不等式。
\end{proof}

\begin{definition}\label{g3:def:measured-gh-convergence}
	一个\term[duliang cedu kongjian]{度量测度空间}[metric measure space]是带有 Borel 测度的度量空间 \((X,d,\nu)\)。本节使用紧致度量概率空间，即 \(X\) 非空紧致且 \(\nu(X)=1\)。若存在紧度量空间 \(Z\) 及等距嵌入 \(\iota_i:X_i\to Z\)、\(\iota:X\to Z\)，使
	\begin{equation}\label{g3:eq:measured-gh-definition}
		d_H^Z(\iota_i(X_i),\iota(X))\longrightarrow0,
		\qquad (\iota_i)_\#\nu_i\rightharpoonup\iota_\#\nu,
	\end{equation}
	就称 \((X_i,d_i,\nu_i)\) \term[daicedu GH shoulian]{带测度 GH 收敛}[measured Gromov--Hausdorff convergence]于 \((X,d,\nu)\)，简称 mGH 收敛。若每个非空开球都有正测度，就称 \(\nu\) 具有\term[man zhicheng cedu]{满支撑}[full support]，即 \(\operatorname{spt}\nu=X\)。
\end{definition}

\begin{proposition}\label{g3:prop:gh-common-compact-ambient}
	设非空紧度量空间 \(X_i\) GH 收敛于 \(X\)，且每个 \(X_i\) 带有 Borel 概率测度 \(\nu_i\)。则它们可等距嵌入同一个紧度量空间，使其 Hausdorff 收敛于 \(X\)；并且存在子列和 \(X\) 上的概率测度 \(\nu\)，使这条子列 mGH 收敛于 \((X,d,\nu)\)。
\end{proposition}
\begin{proof}
	由\cref{g3:thm:gh-correspondence}和桥度量构造，存在 \(\delta_i>0\)、\(\delta_i\to0\)，以及 \(X_i\sqcup X\) 上的度量 \(D_i\)，使两边的包含映射等距，跨边距离至少为 \(\delta_i\)，并有 \(d_H^{D_i}(X_i,X)\le\delta_i\)。在所有这些空间中保留同一份 \(X\)，对 \(a\in X_i,b\in X_j,i\ne j\) 定义
	\[
		D(a,b)=\inf_{x\in X}\bigl(D_i(a,x)+D_j(x,b)\bigr),
	\]
	其余距离沿用各片上的距离。这个公式把不同片通过 \(X\) 连接。任何经过其他片、再回到 \(X\) 的路段，其长度都不小于两个端点在 \(X\) 中的距离；将它替换为该距离不会增长路径长度。因此跨片路径的下确界满足三角不等式，而且不会缩短任何一片原来的距离。不同片之间的距离至少为对应的 \(\delta_i+\delta_j\)，各片到 \(X\) 的距离至少为 \(\delta_i\)，故正定性也成立。

	所得并空间全有界：给定 \(\varepsilon>0\)，所有充分靠后的 \(X_i\) 都落在 \(X\) 的 \(\varepsilon/2\)-邻域中，可由 \(X\) 的有限 \(\varepsilon/2\)-网控制；剩下有限多个紧空间分别取有限网。取完备化即得紧度量空间 \(Z\)。原来的各紧片在其中闭，且仍有 \(d_H^Z(X_i,X)\le\delta_i\)。

	把 \(\nu_i\) 推到 \(Z\)，用\cref{g3:lem:compact-probability-subsequence}取弱收敛子列，极限记为 \(\widetilde\nu\)。由于
	\[
		0\le\int_Zd(z,X)\,\mathrm d\nu_i(z)\le\delta_i,
	\]
	弱收敛给出该连续函数对 \(\widetilde\nu\) 的积分为零。因此 \(\widetilde\nu(Z\setminus X)=0\)，它是 \(X\) 上某个概率测度的推前，这正是所需的 mGH 子列。
\end{proof}

\begin{theorem}[归一化 Ricci 极限的预紧性]\label{g3:thm:ricci-measured-precompactness}
	设连通闭 \(n\) 维黎曼流形序列 \((M_i,g_i)\) 满足式\eqref{g3:eq:ricci-limit-class}，并以式\eqref{g3:eq:ricci-probability-normalization}定义 \(\nu_i\)。则存在 mGH 收敛子列，其极限 \((X,d,\nu)\) 为紧致测地度量概率空间，且 \(\operatorname{spt}\nu=X\)。对任意 \(x\in X\)、\(r>0\)，有
	\begin{equation}\label{g3:eq:ricci-limit-ball-convergence}
		\nu(\partial B_X(x,r))=0,
		\qquad \nu_i(B_{M_i}(x_i,r))\longrightarrow\nu(B_X(x,r))
	\end{equation}
	只要 \(x_i\to x\) 于实现收敛的共同紧致环境中。此外，对 \(0<r<R\)，有
	\begin{equation}\label{g3:eq:ricci-limit-measure-comparison}
		\frac{\nu(B_X(x,r))}{\nu(B_X(x,R))}
		\ge\frac{V_{-K,n}(r)}{V_{-K,n}(R)}.
	\end{equation}
\end{theorem}
\begin{proof}
	\cref{g3:prop:ricci-limit-volume-packing}给出统一直径和覆盖数界；\cref{g3:thm:gh-precompactness}给出 GH 子列，上一命题再给出测度子列。\cref{g3:prop:gh-geodesic-limit}表明 \(X\) 是测地空间。以下将所有空间视为同一紧空间 \(Z\) 的子集。

	固定 \(x\in X\)，取 \(x_i\in M_i\) 使 \(x_i\to x\)。对 \(r>0\) 和 \(0<\eta<r/3\)，开集
	\[
		G_\eta=\{z\in Z:r-\eta<d_Z(x,z)<r+\eta\}
	\]
	包含 \(X\) 中距离 \(x\) 恰为 \(r\) 的点。当 \(i\) 充分大时，它与 \(M_i\) 的交包含于以 \(x_i\) 为中心、内外半径为 \(r-2\eta,r+2\eta\) 的环带。Bishop--Gromov 比较及 \(\nu_i(M_i)=1\) 给出
	\[
		\nu_i(G_\eta)
		\le 1-\frac{V_{-K,n}(r-2\eta)}{V_{-K,n}(r+2\eta)}.
	\]
	用式\eqref{g3:eq:weak-measure-closed-open}的开集不等式，再令 \(\eta\downarrow0\)，得到 \(\nu(\{d(x,\cdot)=r\})=0\)，特别是球边界零测。对固定环境球使用开、闭集不等式，得到其测度收敛；再以半径 \(r\pm\eta\) 的固定球夹住 \(B_{M_i}(x_i,r)\)，令 \(\eta\downarrow0\)，即得式\eqref{g3:eq:ricci-limit-ball-convergence}。

	当 \(0<r<R_0\) 时，式\eqref{g3:eq:ricci-limit-relative-volume}传到极限给出
	\[
		\nu(B_X(x,r))\ge
		\frac{V_{-K,n}(r)}{V_{-K,n}(R_0)}>0.
	\]
	这证明满支撑。最后把流形上的 Bishop--Gromov 球体积比传到极限，分母由满支撑严格为正，得到式\eqref{g3:eq:ricci-limit-measure-comparison}。
\end{proof}

\begin{example}\label{g3:exm:weighted-circle-collapse}
	在二维环面 \((\mathbb R/2\pi\mathbb Z)^2\) 上取坐标 \((x,\theta)\)，令
	\[
		g_\varepsilon=\mathrm dx^2+\varepsilon^2a(x)^2\mathrm d\theta^2,
		\qquad a(x)=2+\sin x,\quad 0<\varepsilon\le1.
	\]
	非零 Christoffel 符号为
	\[
		\Gamma^x_{\theta\theta}=-\varepsilon^2aa',
		\qquad \Gamma^\theta_{x\theta}=
		\Gamma^\theta_{\theta x}=a'/a.
	\]
	由曲率定义计算得 \(K_{g_\varepsilon}=-a''/a=\sin x/(2+\sin x)\ge-1\)，所以 \(\operatorname{Ric}_{g_\varepsilon}\ge-g_\varepsilon\)。投影 \(\pi(x,\theta)=x\) 到长为 \(2\pi\) 的圆周 \(S^1\) 满足
	\[
		d_{S^1}(\pi p,\pi q)\le d_{g_\varepsilon}(p,q)
		\le d_{S^1}(\pi p,\pi q)+3\pi\varepsilon.
	\]
	下界来自投影不增加长度；上界先沿固定 \(\theta\) 的最短底圆弧走，再在终点纤维中走不超过半周。因此 \(\operatorname{diam}(M,g_\varepsilon)\le4\pi\)，且对应 \(\{(p,\pi p)\}\) 的失真至多为 \(3\pi\varepsilon\)，给出 GH 收敛到 \(S^1\)。

	面积元与总体积分别为
	\[
		\mathrm d\operatorname{vol}_{g_\varepsilon}
		=\varepsilon a(x)\,\mathrm dx\,\mathrm d\theta,
		\qquad \operatorname{Vol}_{g_\varepsilon}M=8\pi^2\varepsilon.
	\]
	故未归一化的面积测度总质量趋于零，而归一化后的推前对每个 \(\varepsilon\) 都恰为
	\begin{equation}\label{g3:eq:weighted-collapse-limit-measure}
		\pi_\#\nu_\varepsilon=
		\frac{2+\sin x}{4\pi}\,\mathrm dx.
	\end{equation}
	在桥度量构造中可使对应点 \(p,\pi p\) 的距离趋于零。对共同紧致环境中的连续函数，一致连续性便把 \(\int f(p)\,\mathrm d\nu_\varepsilon\) 与 \(\int f(\pi p)\,\mathrm d\nu_\varepsilon\) 的差控制到零。因此式\eqref{g3:eq:weighted-collapse-limit-measure}确为 mGH 极限测度。它不是圆周的归一化长度测度：在较粗的纤维之上，每单位底圆长度留下更多质量。
\end{example}

\subsection{正则点、切锥及奇异集}
\label{g3:subsec:2903:03}
在光滑流形上，放大一个点附近的距离会得到该点的 Euclid 切空间。度量极限中也可以放大，但不同尺度的子列可能产生不同的极限。定义正则点时，必须说明这种选择是否允许影响结论。

\begin{proposition}\label{g3:prop:ricci-limit-tangents-exist}
	设 \((X,d,\nu)\) 是\cref{g3:thm:ricci-measured-precompactness}给出的极限。则对任意 \(x\in X\) 和任意 \(r_j\downarrow0\)，序列 \((X,r_j^{-1}d,x)\) 都有指向 GH 收敛子列，其极限为固有测地空间。
\end{proposition}
\begin{proof}
	只需验证\cref{g3:thm:pointed-gh-precompactness}的各有界球覆盖数界。去掉有限项，设 \(0<r_j\le1\)。在放大距离 \(d_j=r_j^{-1}d\) 下，式\eqref{g3:eq:ricci-limit-measure-comparison}中的模型曲率变为 \(-Kr_j^2\)，因为
	\[
		V_{-Kr_j^2,n}(t)=r_j^{-n}V_{-K,n}(r_jt).
	\]
	在 \(\overline B_{d_j}(x,R)\) 中取 \(\varepsilon\)-分离点 \(y_1,\ldots,y_N\)。半径 \(\varepsilon/2\) 的开球两两不交，且包含于 \(B_{d_j}(x,R+\varepsilon/2)\)。又有
	\(B_{d_j}(x,R+\varepsilon/2)\subset B_{d_j}(y_a,2R+\varepsilon)\)，所以比较不等式给出
	\[
		\nu(B_{d_j}(y_a,\varepsilon/2))
		\ge\frac{V_{-Kr_j^2,n}(\varepsilon/2)}
		{V_{-Kr_j^2,n}(2R+\varepsilon)}
		\nu(B_{d_j}(x,R+\varepsilon/2)).
	\]
	求和并消去正测度后，得到 \(N\) 不超过上式模型体积比的倒数。对于 \(0\le\kappa\le K\)，模型径向函数满足
	\(t\le s_{-\kappa}(t)\le t e^{\sqrt Kt}\)，从而
	\[
		\frac{V_{-\kappa,n}(2R+\varepsilon)}
		{V_{-\kappa,n}(\varepsilon/2)}
		\le\left(\frac{2R+\varepsilon}{\varepsilon/2}\right)^n
		e^{(n-1)\sqrt K(2R+\varepsilon)}.
	\]
	右端与 \(j\) 无关。取极大分离集即得所需有限覆盖数，指向预紧性定理于是适用。
\end{proof}

\begin{definition}\label{g3:def:ricci-limit-regular-singular}
	设 \((X,d)\) 为固有测地空间，且在每个点的任意趋零放大序列都有指向 GH 收敛子列。沿用\cref{g3:def:metric-tangent-limit}，任何指向 GH 极限
	\[
		(X,r_j^{-1}d,x)\longrightarrow(Y,d_Y,o),\qquad r_j\downarrow0,
	\]
	称为 \(X\) 在 \(x\) 的一个\term[qiezhui]{切锥}[tangent cone]。若对某个非负整数 \(k\)，在 \(x\) 的每一个切锥都以基点保持的方式等距于 \((\mathbb R^k,0)\)，就称 \(x\) 为 \(k\) 维\term[zhengze dian ricci jixian]{正则点}[regular point]，全体此种点记为 \(\mathcal R_k\)。令
	\[
		\mathcal R=\bigcup_{k\ge0}\mathcal R_k,
		\qquad \mathcal S=X\setminus\mathcal R,
	\]
	分别称为正则集和\term[qiyi ji ricci jixian]{奇异集}[singular set]。这里 \(\mathbb R^0\) 表示单点。
\end{definition}

“存在一个 Euclid 切锥”只涉及一列尺度；上面的正则性涉及全部趋零尺度。对固定的 \(k\)，由于上一命题给出所有放大序列的预紧性，“所有切锥都是 \(\mathbb R^k\)”等价于当 \(r\downarrow0\) 时整个族 \((X,r^{-1}d,x)\) 指向 GH 收敛到 \((\mathbb R^k,0)\)。若整族不收敛，就可选出在某个固定半径球上始终偏离 Euclid 模型的序列；取其收敛子列会与所有切锥的规定矛盾。

\ProseParagraphBreak
这里的“切锥”首先是放大极限的名称。真正的\term[duliang zhui]{度量锥}[metric cone] \(C(Z)\) 是将 \([0,\infty)\times Z\) 的 \(\{0\}\times Z\) 压成一点后，赋予余弦距离
\begin{equation}\label{g3:eq:ricci-metric-cone-distance}
	 d_C((r,z),(s,w))^2
	 =r^2+s^2-2rs\cos d_Z(z,w),
	 \qquad \operatorname{diam}Z\le\pi.
\end{equation}
在本节所用情形中，\(Z\) 为非空紧长度空间，这一公式给出锥的长度度量。可以只用有限拼接的路径来验证：每一段的角坐标沿 \(Z\) 中一条单位速度最短弧 \(\eta\) 运动，写成 \(z(t)=\eta(u(t))\)，其中 \(r,u\) 分段光滑，并允许路径经过顶点。给这样的每一段赋予长度
\[
	 L(\gamma)=\int\sqrt{|\dot r|^2+r^2|\dot u|^2}\,\mathrm dt.
\]
给定两端点，沿 \(Z\) 中连接其角坐标的最短弧构成一个角度不超过 \(\pi\) 的 Euclid 扇形；扇形内的直线段给出余弦式的上界。反过来，对任意上述路径取角变量 \(\theta(t)=d_Z(z(0),z(t))\)。距离函数的 \(1\)-Lipschitz 性给出每段上几乎处处的 \(|\dot\theta|\le|\dot u|\)，故平面曲线 \(r(t)e^{\mathrm i\theta(t)}\) 的长度不超过原路径长度。其弦长正是余弦式右端的平方根，因而给出下界；经过顶点的路径长度至少为两端径向距离之和，也满足该下界。这说明公式确为这些路径长度的下确界，并由扇形直线段实现；于是它给出测地度量。径向伸缩是任意正比例的相似变换；切锥定义本身并未要求这些变换存在。

\begin{example}\label{g3:exm:cone-density-singularity}
	取总角度 \(0<\alpha<2\pi\) 的平坦圆锥 \(C_\alpha\)，见\cref{g3:exm:flat-metric-cone}。在顶点 \(o\) 外，度量为 \(\mathrm dr^2+r^2\mathrm d\theta^2\)，其中 \(\theta\) 的周期是 \(\alpha\)。因而
	\[
		\mathcal H^2(B(o,r))
		=\int_0^r\int_0^\alpha t\,\mathrm d\theta\,\mathrm dt
		=\frac\alpha2r^2.
	\]
	径向伸缩表明顶点的所有切锥都等距于 \(C_\alpha\) 本身。它不等距于 \(\mathbb R^2\)，因为后者的半径 \(r\) 球面积为 \(\pi r^2\)；锥的 Hausdorff 维数为二，也排除了其他维数的 Euclid 空间。因此顶点奇异，而顶点外每一点都正则。

	在乘积 \(\mathbb R^{n-2}\times C_\alpha\) 中，奇异集是 \(\mathbb R^{n-2}\times\{o\}\)。在该轴上的点，以 Fubini 积分计算乘积球体积，得到
	\[
		\mathcal H^n(B((0,o),r))
		=\frac\alpha{2\pi}\,\omega_nr^n.
	\]
	理由是对每个 \(z\in\mathbb R^{n-2}\)、\(|z|<r\)，圆锥截片的面积恒为相应 Euclid 二维圆盘面积的 \(\alpha/(2\pi)\) 倍。这一模型把“几乎处处 Euclid”与一条维数为 \(n-2\) 的奇异轴同时呈现出来。
\end{example}

\subsection{所选结构定理的准确范围}
\label{g3:subsec:2903:04}
体积比较已经保证了极限测度的满支撑和切锥的存在。把非塌缩极限的测度识别为 \(\mathcal H^n\)，并控制奇异集的维数，则需要更深的几何分析。下面采用 Cheeger--Colding\footnote{Tobias Holck Colding（1963--），丹麦数学家，研究微分几何、几何分析和偏微分方程，在 Ricci 曲率界和极小曲面理论方面作出贡献。}理论的紧致非塌缩版本；\(\mathcal H^n\) 使用\cref{g3:def:hausdorff-measure}中与 Euclid 体积一致的归一化。

\begin{theorem}[Cheeger--Colding 结构定理]\label{g3:thm:cheeger-colding-noncollapsed}
	设整数 \(n\ge2\)、\(K\ge0\)、\(D,v>0\) 固定，\((M_i^n,g_i)\) 为连通闭黎曼流形，满足
	\[
		\operatorname{Ric}_{g_i}\ge-(n-1)Kg_i,
		\operatorname{diam}(M_i,g_i)\le D,
		\operatorname{Vol}_{g_i}M_i\ge v.
	\]
	若 \((M_i,d_{g_i})\) GH 收敛于 \((X,d)\)，则成立：
	\begin{enumerate}
		\item \(\dim_{\mathrm H}X=n\)。对任意在共同环境中收敛的 \(x_i\to x\) 及任意 \(r>0\)，有
		\begin{equation}\label{g3:eq:noncollapsed-volume-convergence}
			\operatorname{Vol}_{g_i}B(x_i,r)
			\longrightarrow\mathcal H^n(B_X(x,r)).
		\end{equation}
		而且每个归一化体积极限测度都是 \(\mathcal H^n\) 的常数倍。
		\item 每一点的每一个切锥都是式\eqref{g3:eq:ricci-metric-cone-distance}意义下的度量锥，其截面为直径至多 \(\pi\) 的紧长度空间。
		\item \(\mathcal R=\mathcal R_n\)，而且只要某一点存在一个等距于 \(\mathbb R^n\) 的切锥，该点就属于 \(\mathcal R_n\)。
		\item 奇异集满足 \(\dim_{\mathrm H}\mathcal S\le n-2\)。特别地，\(\mathcal H^n(\mathcal S)=0\)。
	\end{enumerate}
\end{theorem}

上述结论分别来自\cite[定理5.9和5.11，第434页；定理5.2，第432--433页；定理6.1，第435--437页]{CheegerColding1997StructureI}。原文采用连通完备无边界流形的指向 GH 极限、\(\operatorname{Ric}\ge-(n-1)g\) 及基点单位球体积下界。为对应这里的条件，先令 \(\widetilde g_i=(K+1)g_i\)，则新的 Ricci 下界满足该归一化；统一直径和总体积下界由\cref{g3:prop:ricci-limit-volume-packing}给出统一的单位球体积下界。最后把距离和体积分别按 \(\sqrt{K+1}\) 和 \((K+1)^{n/2}\) 缩回，即得所述版本。此处外引的是这些结构结论及其证明；前面已经证明的测度子列构造并不包含这一正则性论证。

\begin{corollary}\label{g3:cor:noncollapsed-canonical-measure}
	在\cref{g3:thm:cheeger-colding-noncollapsed}的假设下，归一化黎曼体积沿原来的 GH 收敛序列 mGH 收敛，极限测度唯一且为
	\begin{equation}\label{g3:eq:noncollapsed-canonical-measure}
		\nu=\frac{\mathcal H^n}{\mathcal H^n(X)},
		\qquad v\le\mathcal H^n(X)\le V_{-K,n}(2D+1).
	\end{equation}
	正则集 \(\mathcal R_n\) 在 \(X\) 中稠密，且其补集为 \(\nu\)-零测集。
\end{corollary}
\begin{proof}
	由\cref{g3:prop:gh-diameter}有 \(\operatorname{diam}X\le D\)。取 \(r>D\)，体积收敛式给出
	\(\operatorname{Vol}_{g_i}M_i\to\mathcal H^n(X)\)，从而得到所示上下界。固定一个实现 GH 收敛的共同紧致环境。按上一外部定理第一项，每个概率测度弱极限 \(\mu\) 都是 \(\mathcal H^n\) 的常数倍；总质量为一便唯一确定该常数为 \(\mathcal H^n(X)^{-1}\)。原文以单位球体积归一化与这里以总体积归一化的关系也可直接核对：固定 \(p_i\to p\)，有
	\[
		\frac{\operatorname{Vol}_{g_i}B(p_i,1)}
		{\operatorname{Vol}_{g_i}M_i}
		\longrightarrow
		\frac{\mathcal H^n(B_X(p,1))}{\mathcal H^n(X)}>0.
	\]
	因此两种归一化之间只相差一个有正极限的常数因子。

	任一子列都有进一步的弱收敛子列，而且极限相同；若整列不弱收敛，就会有某个连续测试函数的积分沿一条子列始终与目标相差一个固定正数，矛盾。因此整列 mGH 收敛。最后 \(\mathcal H^n(\mathcal S)=0\) 给出 \(\nu(\mathcal S)=0\)，满支撑保证每个非空开球都有正 \(\nu\) 质量，因而都与 \(\mathcal R_n\) 相交。
\end{proof}

定理中的第三项把“某一个 \(n\) 维 Euclid 切锥”提升为“全部切锥都是该 Euclid 空间”，这是非塌缩结构定理的结论。对于奇异点，第二项只确定切锥属于度量锥这一类，并不确定所有截面相同。第四项则按 Hausdorff 维数控制例外集；它描述几乎处处的度量结构，并未赋予整个极限空间一套光滑黎曼坐标。塌缩时，\cref{g3:exm:weighted-circle-collapse}已经显示，即使距离极限是一条光滑圆周，归一化测度也仍可能带有非恒定密度。
% 纲要条目排为编号小节标题，并进入目录与书签；说明段仍保留于纲要注释中。

% !TeX root = ../../Book3.tex
% 纲要标题：### 29.4 合成曲率与弱微分
% 纲要位置：几何学计划纲要.md，第 3340 行。
% 正文状态：初稿，尚未正式审核。

% 纲要内容开始
% 
% * Cheeger 能量与度量 Sobolev 空间
% * 最优输运和曲率维数条件
% * 光滑流形与合成框架的比较
% * 合成曲率的稳定性与正则性
% 
% ---
% 
% 纲要内容结束

\section{合成曲率与弱微分}
\label{g3:sec:2904}

在几何极限中，距离和测度可以保留下来，坐标中的导数却未必还有意义。要继续讨论能量与曲率，须把它们改写成只使用距离、测度及极限的条件。本节采用两种改写：以 Lipschitz 函数的局部斜率松弛出一阶能量，以概率测度的最优输运表达 Ricci 下界。它们分别保留微分的大小与曲率的凸性作用，合在一起便得到一种适用于奇异空间的黎曼结构。

\subsection{Cheeger 能量与度量 Sobolev 空间}
\label{g3:subsec:2904:01}
先设 \((X,d)\) 为紧致度量空间，\(\mathfrak m\) 为满支集的 Borel 概率测度。满支集是指每个非空开球都有正测度。以下函数取实值，\(L^2(\mathfrak m)\) 按几乎处处相等识别函数。对 Lipschitz 函数 \(f\)，定义它在 \(x\) 处的\term[jubu xielv]{局部斜率}[local slope]
\begin{equation}\label{g3:eq:metric-local-slope}
	\operatorname{lip}f(x)\coloneqq
	\limsup_{\substack{y\to x\\y\ne x}}
	\frac{|f(y)-f(x)|}{d(x,y)}.
\end{equation}
孤立点处约定该值为零。斜率是 Borel 函数，且不超过 \(f\) 的全局 Lipschitz 常数；可用一组可数稠密点计算各个邻域内的上确界来核对可测性。在光滑流形上，这个量在可微点恰等于 \(|\nabla f|\)。

\ProseParagraphBreak
直接积分斜率平方，只能处理已经具有 Lipschitz 代表的函数。与第21章用光滑逼近研究变差的思路相同，可以把允许的函数扩展到 \(L^2\)，并把极限中能够保留的最小能量作为定义。

\begin{definition}\label{g3:def:cheeger-energy}
	设 \((X,d,\mathfrak m)\) 为紧致度量概率空间，\(\operatorname{supp}\mathfrak m=X\)。记 \(\operatorname{lip}v(x)=\limsup_{y\to x,y\ne x}|v(y)-v(x)|/d(x,y)\)，孤立点处取零。其\term[Cheeger nengliang]{Cheeger 能量}[Cheeger energy]是
	\begin{equation}\label{g3:eq:cheeger-relaxation}
		\operatorname{Ch}(f)\coloneqq
		\inf_{\substack{f_j\in\operatorname{Lip}(X)\\f_j\to f\text{ 于 }L^2(\mathfrak m)}}
		\liminf_{j\to\infty}\frac12\int_X(\operatorname{lip}f_j)^2\,\mathrm d\mathfrak m,
		\qquad f\in L^2(\mathfrak m).
	\end{equation}
	取不到有限值时，令 \(\operatorname{Ch}(f)=+\infty\)。相应的\term[duliang Sobolev kongjian]{度量 Sobolev 空间}[metric Sobolev space]为
	\[
		W^{1,2}(X,d,\mathfrak m)\coloneqq
		\{f\in L^2(\mathfrak m):\operatorname{Ch}(f)<\infty\},
		\qquad
		\|f\|_{W^{1,2}}^2\coloneqq\|f\|_2^2+2\operatorname{Ch}(f).
	\]
\end{definition}
\symidx{\(\operatorname{lip}v\)：函数的局部 Lipschitz 斜率}
\symidx{\(\operatorname{Ch}\)：度量测度空间的松弛 Cheeger 能量}

这称为原斜率能量关于 \(L^2\) 收敛的\term[songchi nengliang]{松弛能量}[relaxed energy]。定义中的 \(1/2\) 与前面各章的 Dirichlet 能量约定一致。紧致度量空间上的连续函数可由 Lipschitz 函数一致逼近，例如对连续 \(h\) 使用 \(h_j(x)=\inf_y\{h(y)+jd(x,y)\}\)。连续函数又在 \(L^2(\mathfrak m)\) 中稠密，因此定义中的逼近序列总能找到；困难在于其能量是否有界。

\begin{proposition}\label{g3:prop:cheeger-basic-properties}
	设 \((X,d,\mathfrak m)\) 为满支集紧致度量概率空间，\(\operatorname{Ch}\) 及 \(W^{1,2}\) 分别为\cref{g3:def:cheeger-energy}定义的松弛斜率能量及其有限能量空间。则\(\operatorname{Ch}\) 非负、凸，在 \(L^2(\mathfrak m)\) 上下半连续，且
	\[
		\operatorname{Ch}(af)=a^2\operatorname{Ch}(f),\qquad
		\sqrt{\operatorname{Ch}(f+g)}
		\le\sqrt{\operatorname{Ch}(f)}+\sqrt{\operatorname{Ch}(g)}
	\]
	对有限能量函数及实数 \(a\) 成立。上述 \(W^{1,2}\) 是 Banach 空间。
\end{proposition}
\begin{proof}
	局部斜率满足 \(\operatorname{lip}(af)=|a|\operatorname{lip}f\) 及
	\(\operatorname{lip}(f+g)\le\operatorname{lip}f+\operatorname{lip}g\)。若 \(\operatorname{Ch}(f)<\infty\)，从定义中逐次选取逼近序列及其足够靠后的项，便得到 \(f_j\to f\)，使斜率能量趋向 \(\operatorname{Ch}(f)\)。分别对 \(f,g\) 这样选取，再使用积分的 Minkowski 不等式，得到所述三角不等式。缩放逼近序列并反向缩放，得到齐次性；\(a=0\) 时直接用零函数。能量平方根的三角不等式与齐次性使它成为半范数，进而其平方 \(\operatorname{Ch}\) 为凸泛函。

	证明下半连续性只需考虑 \(f_j\to f\) 且 \(\liminf_j\operatorname{Ch}(f_j)<\infty\)。先取实现下极限的子列，再选 Lipschitz 函数 \(v_j\)，使
	\[
		\|v_j-f_j\|_2<j^{-1},\qquad
		\frac12\int_X(\operatorname{lip}v_j)^2\,\mathrm d\mathfrak m
		\le\operatorname{Ch}(f_j)+j^{-1}.
	\]
	于是 \(v_j\to f\)，定义直接给出 \(\operatorname{Ch}(f)\le\liminf_j\operatorname{Ch}(f_j)\)。

	两个半范数 \(\|f\|_2\)、\(\sqrt{2\operatorname{Ch}(f)}\) 的平方和定义一个范数。若 \((f_j)\) 在此范数下为 Cauchy 列，先取其 \(L^2\) 极限 \(f\)。下半连续性给出
	\[
		\operatorname{Ch}(f_j-f)
		\le\liminf_{k\to\infty}\operatorname{Ch}(f_j-f_k).
	\]
	右端随 \(j\to\infty\) 趋于零。再用半范数三角不等式知 \(\operatorname{Ch}(f)<\infty\)，故 \(f_j\to f\) 于 \(W^{1,2}\)。
\end{proof}

二次齐次性还不足以使能量来自内积。称 Cheeger 能量为\term[erci Cheeger nengliang]{二次 Cheeger 能量}[quadratic Cheeger energy]，是指它在有限能量定义域上满足平行四边形恒等式
\begin{equation}\label{g3:eq:cheeger-quadratic}
	\operatorname{Ch}(f+g)+\operatorname{Ch}(f-g)
	=2\operatorname{Ch}(f)+2\operatorname{Ch}(g).
\end{equation}
此时极化公式
\begin{equation}\label{g3:eq:cheeger-polarization}
	\mathcal E(f,g)\coloneqq
	\frac12\bigl[\operatorname{Ch}(f+g)-\operatorname{Ch}(f-g)\bigr]
\end{equation}
给出对称双线性型，且 \(\mathcal E(f,f)=2\operatorname{Ch}(f)\)。于是 \(W^{1,2}\) 的范数由内积 \((f,g)_{L^2}+\mathcal E(f,g)\) 诱导。反过来，若该范数来自内积，减去 \(L^2\) 范数的平行四边形恒等式便得到式\eqref{g3:eq:cheeger-quadratic}。具有这一性质的度量测度空间称为\term[wuqiongxiao Hilbert kongjian]{无穷小 Hilbert 空间}[infinitesimally Hilbertian space]。

\begin{example}\label{g3:exm:nonquadratic-norm-energy}
	令 \(X=[0,1]^2\)，\(d(x,y)=\max\{|x_1-y_1|,|x_2-y_2|\}\)，\(\mathfrak m\) 为二维 Lebesgue 测度。若 \(u(x)=ax_1+bx_2\)，则
	\begin{equation}\label{g3:eq:sup-norm-affine-energy}
		\operatorname{Ch}(u)=\frac12(|a|+|b|)^2.
	\end{equation}
	先在正方形内部沿各方向计算，得 \(\operatorname{lip}u=|a|+|b|\)；边界为零测集，所以常值逼近序列给出上界。

	为证明下界，取任意能量有界的 Lipschitz 逼近 \(u_j\to u\)。两种坐标范数等价，\cref{g3:thm:rademacher}使 \(u_j\) 几乎处处可微，其经典导数也是弱导数。取 \(v=(\operatorname{sgn}a,\operatorname{sgn}b)\)，在 \((a,b)\ne(0,0)\) 时有 \(\|v\|_\infty=1\)，故
	\[
		|\mathrm D_vu_j|\le\operatorname{lip}u_j\quad\text{几乎处处}.
	\]
	方向导数在 \(L^2\) 中有界。其弱收敛子列的极限由分布配对确定为 \(\mathrm D_vu=|a|+|b|\)，范数的弱下半连续性因此给出式\eqref{g3:eq:sup-norm-affine-energy}的下界；\(a=b=0\) 时结论直接成立。

	现在取 \(f=x_1\)、\(g=x_2\)。则 \(\operatorname{Ch}(f)=\operatorname{Ch}(g)=1/2\)，而 \(\operatorname{Ch}(f+g)=\operatorname{Ch}(f-g)=2\)。式\eqref{g3:eq:cheeger-quadratic}的两边分别为 \(4\) 和 \(2\)，所以该空间不具有无穷小 Hilbert 性。距离由一个范数产生，并不能保证导数能量是二次型。
\end{example}

\subsection{最优输运和曲率维数条件}
\label{g3:subsec:2904:02}
能量测量单个函数的变化，Ricci 下界则可以通过测度之间的移动表达。以下进一步假定 \((X,d)\) 为紧致测地空间，并仍取满支集概率测度 \(\mathfrak m\)。记 \(\mathcal P(X)\) 为 \(X\) 上的 Borel 概率测度全体；由于 \(X\) 紧致，所有这些测度都有有限二阶矩。

\begin{definition}\label{g3:def:quadratic-transport-distance}
	设 \((X,d)\) 为非空紧致度量空间，\(\mathcal P(X)\) 为其 Borel 概率测度空间。
	给定 \(\nu_0,\nu_1\in\mathcal P(X)\)，一个\term[ouhe]{耦合}[coupling]是 \(X\times X\) 上的概率测度 \(\pi\)，使两个坐标投影的推前分别等于 \(\nu_0,\nu_1\)。记所有耦合组成的集合为 \(\Pi(\nu_0,\nu_1)\)。定义\term[erci shuyun juli]{二次输运距离}[quadratic Wasserstein distance]\footnote{\OriginalName{Леонид Нисонович Васерштейн}（Leonid Nisonovich Vaserstein，1944--），出生于苏联、长期在美国工作的数学家，研究代数与动力系统；人物资料见\href{https://www.mathnet.ru/php/person.phtml?option_lang=eng\&personid=23042}{Math-Net 人物档案}及\href{https://science.psu.edu/math/people/lxv1}{Penn State 教师主页}。}
	\begin{equation}\label{g3:eq:wasserstein-two}
		W_2(\nu_0,\nu_1)^2\coloneqq
		\inf_{\pi\in\Pi(\nu_0,\nu_1)}
		\int_{X\times X}d(x,y)^2\,\mathrm d\pi(x,y).
	\end{equation}
	达到下确界的耦合称为\term[zuiyou shuyun jihua]{最优输运计划}[optimal transport plan]。若概率测度曲线 \((\nu_t)_{0\le t\le1}\) 满足
	\[
		W_2(\nu_s,\nu_t)=|s-t|W_2(\nu_0,\nu_1)
		\qquad(0\le s,t\le1),
	\]
	则称它为\term[Wasserstein changsu cedixian]{Wasserstein 常速测地线}[constant-speed Wasserstein geodesic]。
\end{definition}
\symidx{\(\mathcal P(X),\Pi(\nu_0,\nu_1)\)：Borel 概率测度空间与指定边缘的耦合集}
\symidx{\(W_2\)：概率测度之间的二次输运距离}

耦合允许把同一位置的质量分送到不同位置，因而不要求输运来自一个单值映射。\(\nu_0\otimes\nu_1\) 是一个耦合，保证可行集非空。紧致性还保证最优计划存在：\(\mathcal P(X\times X)\) 弱紧，边缘条件在弱收敛下封闭，而代价函数 \(d^2\) 连续。这里概率测度的弱紧性可由 \(C(X\times X)\) 的可分性、对可数稠密测试族取对角子列及测度表示定理得到。

\begin{proposition}\label{g3:prop:wasserstein-metric}
	对任意紧致度量空间 \((X,d)\)，式\eqref{g3:eq:wasserstein-two}在 \(\mathcal P(X)\) 上定义距离，且 \(W_2(\delta_x,\delta_y)=d(x,y)\)。
\end{proposition}
\begin{proof}
	非负性和对称性来自代价函数。若最小代价为零，最优计划集中在对角线上，两个边缘便相同。反之，对角推前给出同一测度到自身的零代价计划。两个点质量只有一个耦合，得到末式。

	三角不等式需要把两次输运接起来。设 \(\pi_{01},\pi_{12}\) 是从 \(\nu_0\) 到 \(\nu_1\)、从 \(\nu_1\) 到 \(\nu_2\) 的最优计划。取有限 Borel 分割 \(\{A_i\}\)，使各片直径不超过 \(\varepsilon\)。将 \(\pi_{01}\) 限制在第二坐标属于 \(A_i\) 的部分，记为 \(\alpha_i\)；将 \(\pi_{12}\) 限制在第一坐标属于 \(A_i\) 的部分，记为 \(\beta_i\)。它们的质量均为 \(m_i=\nu_1(A_i)\)。在坐标为 \((x,y,z,w)\) 的 \(X^4\) 上令
	\[
		\gamma_\varepsilon=\sum_{m_i>0}m_i^{-1}\alpha_i\otimes\beta_i.
	\]
	其 \((x,y)\)、\((z,w)\) 边缘分别是原来的两个计划，且 \(d(y,z)\le\varepsilon\) 几乎处处。令 \(\varepsilon\downarrow0\) 并取弱收敛子列，得到集中在 \(y=z\) 上的极限。因此存在 \(X^3\) 上的概率测度 \(\gamma\)，其相邻两个边缘就是 \(\pi_{01},\pi_{12}\)。对它使用距离三角不等式及 Minkowski 不等式，便有
	\[
		\begin{aligned}
		W_2(\nu_0,\nu_2)
		&\le\left(\int d(x,w)^2\,\mathrm d\gamma\right)^{1/2}\\
		&\le\left(\int d(x,y)^2\,\mathrm d\gamma\right)^{1/2}
		+\left(\int d(y,w)^2\,\mathrm d\gamma\right)^{1/2}\\
		&=W_2(\nu_0,\nu_1)+W_2(\nu_1,\nu_2).
		\end{aligned}
	\]
\end{proof}

\begin{example}\label{g3:exm:transport-not-mixture}
	在 \([0,1]\) 上，\(\nu_t=\delta_t\) 是连接 \(\delta_0\) 与 \(\delta_1\) 的常速输运测地线。线性混合 \(\eta_t=(1-t)\delta_0+t\delta_1\) 则满足
	\[
		W_2(\eta_s,\eta_t)^2=|t-s|.
	\]
	因为恰有 \(|t-s|\) 的质量必须在距离为 \(1\) 的两个端点间移动，余下质量可留在原地。因此线性混合不是同一条输运测地线：输运凸性沿质量的空间移动计算。
\end{example}

对参考概率测度 \(\mathfrak m\)，定义\term[xiangdui shang]{相对熵}[relative entropy]
\begin{equation}\label{g3:eq:metric-relative-entropy}
	\operatorname{Ent}_{\mathfrak m}(\nu)\coloneqq
	\begin{cases}
		\displaystyle\int_X\rho\log\rho\,\mathrm d\mathfrak m,
		&\nu=\rho\mathfrak m,\\
		+\infty,&\nu\not\ll\mathfrak m.
	\end{cases}
\end{equation}
约定 \(0\log0=0\)。因为 \(r\log r\ge-1/e\)，该积分有确定的扩展实数值；又由凸性和 \(\int\rho\,\mathrm d\mathfrak m=1\)，其值非负。有限熵测度构成的集合记为 \(D(\operatorname{Ent}_{\mathfrak m})\)。

\begin{definition}\label{g3:def:cd-rcd-infinity}
	设 \((X,d,\mathfrak m)\) 为满支集的紧致测地度量概率空间，\(K\in\mathbb R\)。记 \(\operatorname{Ent}_{\mathfrak m}(\rho\mathfrak m)=\int\rho\log\rho\,\mathrm d\mathfrak m\)，对非绝对连续概率测度取 \(+\infty\)，\(D(\operatorname{Ent}_{\mathfrak m})\) 为有限熵概率测度集，\(W_2\) 为二次输运距离；无穷小 Hilbert 性指 Cheeger 能量满足平行四边形恒等式。若任意 \(\nu_0,\nu_1\in D(\operatorname{Ent}_{\mathfrak m})\) 之间都存在一条 \(W_2\) 常速测地线 \((\nu_t)\)，使
	\begin{equation}\label{g3:eq:cd-infinity-entropy-convexity}
		\operatorname{Ent}_{\mathfrak m}(\nu_t)
		\le(1-t)\operatorname{Ent}_{\mathfrak m}(\nu_0)
		+t\operatorname{Ent}_{\mathfrak m}(\nu_1)
		-\frac K2t(1-t)W_2(\nu_0,\nu_1)^2
	\end{equation}
	对全部 \(t\in[0,1]\) 成立，就称该空间满足\term[qulv weishu tiaojian]{曲率维数条件}[curvature-dimension condition] \(\mathrm{CD}(K,\infty)\)。若它还具有无穷小 Hilbert 性，就称其满足\term[Limanni qulv weishu tiaojian]{黎曼曲率维数条件}[Riemannian curvature-dimension condition] \(\mathrm{RCD}(K,\infty)\)。
\end{definition}
\symidx{\(\operatorname{Ent}_{\mathfrak m}\)：相对于参考测度的 Boltzmann 熵}

\(\mathrm{CD}(K,\infty)\) 要求为每对端点找到一条满足熵不等式的测地线，并未在定义中要求所有候选测地线都满足它。这里采用“普通 \(\mathrm{CD}\) 条件加二次能量”的 \(\mathrm{RCD}\) 定义；其与线性热流等刻画的等价性见 \cite[定理6.1，接受稿第30页]{AmbrosioGigliMondinoRajala2015RCD}。记号中的 \(\infty\) 表示没有施加一个有限的维数上界。有限参数 \(N\) 的 \(\mathrm{CD}(K,N)\)、\(\mathrm{RCD}(K,N)\) 还包含维数信息，不能把本节的条件直接缩写或改记为那些条件。

\subsection{光滑流形与合成框架的比较}
\label{g3:subsec:2904:03}
先核对能量是否恢复了原来的微分。设 \((M,g)\) 为闭连通光滑黎曼流形，\(V\in C^\infty(M)\)，并令
\begin{equation}\label{g3:eq:weighted-manifold-probability}
	\mathrm d\mathfrak m=Z^{-1}e^{-V}\mu_g,
	\qquad Z=\int_Me^{-V}\mu_g.
\end{equation}
这里 \(d_g\) 是黎曼距离，\(\mu_g\) 是体积密度。因 \(M\) 紧致且权重光滑、严格为正，加权和无权的 \(L^2\)、\(W^{1,2}\) 范数互相控制。

\begin{theorem}\label{g3:thm:smooth-cheeger-identification}
	设 \((M^n,g)\) 为闭连通光滑黎曼流形，\(n\ge1\)，\(V\in C^\infty(M)\)，\(\mathrm d\mathfrak m=Z^{-1}e^{-V}\mathrm d\mu_g\)，\(Z=\int_Me^{-V}\mathrm d\mu_g\)。则
	\[
		W^{1,2}(M,d_g,\mathfrak m)=W^{1,2}(M,g),
	\]
	并且对该空间中的每个 \(f\)，有
	\begin{equation}\label{g3:eq:smooth-cheeger-energy}
		\operatorname{Ch}(f)=\frac12\int_M|\nabla_gf|^2\,\mathrm d\mathfrak m.
	\end{equation}
	不属于经典 \(W^{1,2}\) 的 \(L^2\) 函数具有无穷 Cheeger 能量。因此 \((M,d_g,\mathfrak m)\) 具有无穷小 Hilbert 性。
\end{theorem}
\begin{proof}
	先比较局部斜率与微分。在 \(f\) 可微的一点 \(x\)，法坐标给出
	\[
		f(\exp_xv)-f(x)=\mathrm df_x(v)+o(|v|),\qquad
		d_g(x,\exp_xv)=|v|
	\]
	对足够小的 \(v\) 成立。取方向上确界即得 \(\operatorname{lip}f(x)=|\mathrm df_x|_g\)。对 Lipschitz 函数，局部坐标中的 Rademacher 定理保证几乎处处可微，而其微分由沿坐标线的分部积分识别为弱微分；因此该等式几乎处处成立。

	若 \(f\in W^{1,2}(M,g)\)，由\cref{g3:thm:sobolev-smooth-density}可取光滑 \(f_j\to f\) 于经典 \(W^{1,2}\)。闭流形上的光滑函数全局 Lipschitz，权重的上下界又保证加权梯度范数收敛。所以把 \(f_j\) 代入式\eqref{g3:eq:cheeger-relaxation}，得到
	\[
		\operatorname{Ch}(f)\le\frac12\int_M|\nabla f|^2\,\mathrm d\mathfrak m.
	\]
	这也说明每个经典 Sobolev 函数都具有有限的松弛能量。

	反过来，设 \(f_j\to f\) 于 \(L^2(\mathfrak m)\)，且斜率能量的下极限有限。取实现该下极限的子列，便有 \(\nabla f_j\) 在加权 \(L^2(\mathrm TM)\) 中有界，故可再取弱收敛子列 \(\nabla f_j\rightharpoonup G\)。加权与无权范数等价，\(f_j\to f\) 也在无权 \(L^2\) 中成立。对任意光滑向量场 \(Y\)，弱导数恒等式
	\[
		\int_M\langle\nabla f_j,Y\rangle\mu_g
		=-\int_M f_j\operatorname{div}_gY\,\mu_g
	\]
	可传到极限：左端在加权空间中用光滑测试场 \(Ze^VY\) 配对即可。因此 \(G\) 是 \(f\) 的弱梯度，\(f\in W^{1,2}(M,g)\)。Hilbert 范数的弱下半连续性给出
	\[
		\frac12\int_M|\nabla f|^2\,\mathrm d\mathfrak m
		\le\liminf_j\frac12\int_M|\nabla f_j|^2\,\mathrm d\mathfrak m
		=\liminf_j\frac12\int_M(\operatorname{lip}f_j)^2\,\mathrm d\mathfrak m.
	\]
	对全部逼近序列取下确界，便得到反向不等式；同时，有限 Cheeger 能量强制 \(f\) 属于经典 Sobolev 空间。最后，点态内积的平行四边形恒等式在积分后给出式\eqref{g3:eq:cheeger-quadratic}。
\end{proof}

度量确定斜率，测度则决定能量如何加权。同样，合成曲率在加权模型中比较的是
\begin{equation}\label{g3:eq:weighted-ricci-infinity}
	\operatorname{Ric}_{V,\infty}\coloneqq
	\operatorname{Ric}_g+\operatorname{Hess}_gV,
\end{equation}
称为\term[jiaquan Ricci zhangliang]{加权 Ricci 张量}[weighted Ricci tensor]。归一化常数 \(Z\) 对它没有影响，因为把 \(V\) 加上常数不改变 Hessian。

\begin{theorem}\label{g3:thm:smooth-cd-rcd-comparison}
	设 \((M^n,g)\) 为闭连通光滑黎曼流形，\(n\ge1\)，\(V\in C^\infty(M)\)，\(\mathrm d\mathfrak m=Z^{-1}e^{-V}\mathrm d\mu_g\)，\(Z=\int_Me^{-V}\mathrm d\mu_g\)，\(K\in\mathbb R\)。下列条件等价：
	\begin{enumerate}
		\item \(\operatorname{Ric}_g+\operatorname{Hess}_gV\ge Kg\)；
		\item \((M,d_g,\mathfrak m)\) 满足 \(\mathrm{CD}(K,\infty)\)；
		\item \((M,d_g,\mathfrak m)\) 满足 \(\mathrm{RCD}(K,\infty)\)。
	\end{enumerate}
\end{theorem}

第一项与第二项的等价性作为外部定理使用，精确来源为 \cite[定理7.3(b)，第948--949页]{LottVillani2009Transport}；其中第五项就是相对熵的弱 \(K\) 位移凸性。该定理适用于这里的闭光滑加权模型。第二项与第三项的等价性则由\cref{g3:thm:smooth-cheeger-identification}及定义立即得到。因而在无权情形，合成条件恢复通常的 Ricci 下界；在有权情形，必须同时计算势函数的 Hessian。

\begin{example}\label{g3:exm:weighted-circle-synthetic-curvature}
	在单位圆周 \(\mathbb R/(2\pi\mathbb Z)\) 上，给定 \(a\in\mathbb R\)，令 \(Z=\int_0^{2\pi}e^{-a\cos\theta}\,\mathrm d\theta\)，取 \(g=\mathrm d\theta^2\)、\(V(\theta)=a\cos\theta\)。一维 Ricci 张量为零，而
	\[
		\operatorname{Ric}_{V,\infty}=-a\cos\theta\,g,
		\qquad
		\operatorname{Ch}(f)=\frac1{2Z}\int_0^{2\pi}|f'(\theta)|^2e^{-a\cos\theta}\,\mathrm d\theta.
	\]
	因此该加权圆周满足 \(\mathrm{RCD}(-|a|,\infty)\)；当 \(a\ne0\) 时不满足 \(\mathrm{CD}(0,\infty)\)。改变测度而保持圆周距离不变，已经足以改变合成曲率下界。取 \(a=0\) 则恢复平坦圆周的 \(\mathrm{RCD}(0,\infty)\) 条件。
\end{example}

\subsection{合成曲率的稳定性与正则性}
\label{g3:subsec:2904:04}
现在使用\cref{g3:def:measured-gh-convergence}中的紧致测度 Gromov--Hausdorff 收敛：在共同紧致空间中等距嵌入，既要求像集 Hausdorff 收敛，也要求推前概率测度弱收敛。空间的点和距离必须比较，作为能量积分依据的测度也必须比较。

\begin{theorem}\label{g3:thm:compact-rcd-stability}
	设 \(K\in\mathbb R\)，\((X_j,d_j,\mathfrak m_j)\) 为满支集的紧致测地度量概率空间，并在测度 Gromov--Hausdorff 意义下收敛到同类空间 \((X,d,\mathfrak m)\)。若每个 \(X_j\) 都满足 \(\mathrm{RCD}(K,\infty)\)，则极限也满足 \(\mathrm{RCD}(K,\infty)\)。
\end{theorem}

此处采用 \cite[定理7.2，2013年11月19日作者稿第54页]{GigliMondinoSavare2015Stability}的紧致概率特例。该文使用更一般的带基点测度收敛；在本节的共同紧致环境中，选取收敛的基点后，测度弱收敛直接满足其所需条件。极限具有满支集作为本定理的明确假设保留；若测度只支撑在较小闭集上，合成条件应在该支集上讨论。

\ProseParagraphBreak
稳定性同时保留熵的输运凸性与能量的二次性。两者都涉及随空间变化的函数或测度，不能只把一个固定坐标公式逐项取极限。特别是 Cheeger 能量先经过松弛，二次性的传递需要比较这些松弛能量的极限；它并不来自式\eqref{g3:eq:cheeger-quadratic}的形式相同。

\begin{example}\label{g3:exm:collapsed-torus-cheeger-energy}
	取二维环面 \(M_\varepsilon=(\mathbb R/2\pi\mathbb Z)^2\)，度量与归一化体积为
	\[
		g_\varepsilon=\mathrm d\theta^2+\varepsilon^2\mathrm d\varphi^2,
		\qquad
		\mathrm d\mathfrak m_\varepsilon=\frac{\mathrm d\theta\,\mathrm d\varphi}{4\pi^2}.
	\]
	投影到第一圆周使距离的误差至多为 \(\pi\varepsilon\)，且把 \(\mathfrak m_\varepsilon\) 精确推为 \(\mathrm d\theta/(2\pi)\)。因此该族在测度 Gromov--Hausdorff 意义下塌缩到均匀测度圆周。每个度量平坦，故由光滑比较和稳定性，极限满足 \(\mathrm{RCD}(0,\infty)\)。

	能量直接显示消失方向的代价：对光滑函数，
	\[
		\operatorname{Ch}_\varepsilon(f)
		=\frac1{8\pi^2}\int_0^{2\pi}\int_0^{2\pi}
		\left(|\partial_\theta f|^2+\varepsilon^{-2}|\partial_\varphi f|^2\right)
		\,\mathrm d\theta\,\mathrm d\varphi.
	\]
	特别地，\(\operatorname{Ch}_\varepsilon(\cos\theta)=1/4\)，而 \(\operatorname{Ch}_\varepsilon(\cos\varphi)=1/(4\varepsilon^2)\)。沿保留方向变化的能量保持有限，沿收缩方向的这个固定振荡则越来越昂贵。归一化测度保证极限仍有总质量 \(1\)，不会随未归一化面积一起趋于零。
\end{example}

弱微分与合成曲率给极限空间提供了可用的分析结构，尚未给出一套光滑坐标。进一步的正则性问题包括：哪些点有 Euclid 切空间，能否由可数个坐标片覆盖满测度部分，以及不同维数的部分怎样相互邻接。回答这些问题需要区分有限维曲率条件、非塌缩与额外正则性输入；本节的 \(\mathrm{RCD}(K,\infty)\) 稳定性本身没有指定有限维数，也不将奇异极限变成光滑流形。
% 纲要条目排为编号小节标题，并进入目录与书签；说明段仍保留于纲要注释中。


\begin{chaptersummary}
	光滑紧性由局部坐标控制走向整体流形。径向平行标架把法坐标中的联络与余标架写成积分方程，曲率的协变导数界随之转成坐标导数界。有限图册嵌入固定 Euclid 空间后，坐标块排除极限中的额外分支，局部度量相容地粘合，管状投影给出全局比较微分同胚。这样得到的光滑收敛蕴含 GH 收敛和体积收敛。

	\ProseParagraphBreak
	圆丛塌缩保留了水平距离，却压缩了垂直方向的长度和体积；其曲率可以一致有界，注入半径则趋于零。非平凡 Nil 圆丛说明，局部纤维在整体上可以扭转。一般有界曲率塌缩到光滑基底的纤维化结论使用 Fukaya 定理；观察尺度与维数的关系可以在环面上直接计算。

	\ProseParagraphBreak
	Ricci 下界与直径界通过体积比较控制覆盖数，归一化体积则在共同紧环境中取得非零测度极限。非塌缩条件进一步使极限测度由相同维数的 Hausdorff 测度确定，切锥和奇异集受到 Cheeger--Colding 结构定理的约束。距离收敛、测度选择和切锥正则性在这一过程中各有明确的作用。

	\ProseParagraphBreak
	Cheeger 能量以局部斜率的松弛恢复弱微分，光滑加权流形上等于通常的 Dirichlet 能量。概率测度的二次输运距离使相对熵的凸性成为曲率条件；RCD 还要求能量具有二次性。光滑比较把这些条件与加权 Ricci 张量联系起来，稳定性则使同一合成条件可以保留到带测度极限中。
\end{chaptersummary}

\BookExercises

\begin{exercise}\label{g3:ex:metric-comparison-energy}
	设 \(g,h\) 为同一个闭 \(n\) 维流形上的光滑黎曼度量，且 \(e^{-2\delta}g\le h\le e^{2\delta}g\)，其中 \(\delta\ge0\)。对光滑实函数 \(u\) 定义 \(E_g(u)=\int_M|\mathrm du|_g^2\mu_g\)。证明
	\[
		e^{-n\delta}\mu_g\le\mu_h\le e^{n\delta}\mu_g,
		\qquad e^{-(n+2)\delta}E_g(u)\le E_h(u)\le e^{(n+2)\delta}E_g(u).
	\]
	并给出两个距离空间的 GH 距离上界。
\end{exercise}
\begin{solution}
	在 \(g\) 正交标架中，\(h\) 的全部特征值介于 \(e^{-2\delta}\) 与 \(e^{2\delta}\) 之间。行列式平方根给出体积密度界，逆矩阵给出 \(e^{-2\delta}|\mathrm du|_g^2\le|\mathrm du|_h^2\le e^{2\delta}|\mathrm du|_g^2\)。逐点相乘再积分即得能量界。对道路长度积分又有 \(e^{-\delta}d_g\le d_h\le e^\delta d_g\)，所以恒等对应的畸变至多为 \((e^\delta-1)\operatorname{diam}_gM\)。对应公式给
	\[
		d_{\mathrm{GH}}((M,d_g),(M,d_h))
		\le\tfrac12(e^\delta-1)\operatorname{diam}_gM.
	\]
\end{solution}

\begin{exercise}\label{g3:ex:curvature-derivative-control}
	在周期 \(2\pi\) 的平坦二维环面上取 \(u_j=j^{-2}\sin(jx)\)、\(g_j=e^{2u_j}g_0\)。证明 \(g_j\to g_0\) 于固定坐标的 \(C^1\)，曲率一致有界，但 \(\sup|\nabla K_{g_j}|_{g_j}\to\infty\)。由此排除任意子列的 \(C^3\) Cheeger--Gromov 收敛。
\end{exercise}
\begin{solution}
	\(u_j\to0\) 一致且 \(\partial_xu_j=j^{-1}\cos(jx)\to0\) 一致，故度量 \(C^1\) 收敛。二维共形曲率公式给
	\[
		K_{g_j}=e^{-2u_j}\sin(jx),\qquad |K_{g_j}|\le e^2.
	\]
	在 \(x=0\) 处，\(u_j=0\)，\(\partial_xK_{g_j}=j\)，因此 \(|\nabla K_{g_j}|_{g_j}=j\)。若有 \(C^3\) Cheeger--Gromov 收敛子列，拉回后的正定度量及其前三阶导数在紧流形上有统一界；曲率梯度由这些导数及逆度量组成，范数便应一致有界。该范数在微分同胚拉回下不变，与上述发散矛盾。
\end{solution}

\begin{exercise}\label{g3:ex:circle-bundle-basic-energy}
	对\cref{g3:prop:circle-bundle-collapse}中的联络度量，用 \(\nu_\varepsilon=\mu_{g_\varepsilon}/\operatorname{Vol}_{g_\varepsilon}P\) 归一化体积。设 \(v\in C^\infty(B)\)，\(u=v\circ\pi\)。证明
	\[
		\int_Pu^2\,\mathrm d\nu_\varepsilon
		=\frac{1}{\operatorname{Vol}_hB}\int_Bv^2\mu_h,
		\qquad
		\int_P|\mathrm du|_{g_\varepsilon}^2\,\mathrm d\nu_\varepsilon
		=\frac{1}{\operatorname{Vol}_hB}\int_B|\mathrm dv|_h^2\mu_h.
	\]
\end{exercise}
\begin{solution}
	函数 \(u\) 沿纤维常值，所以 \(\mathrm du(V)=0\)，其梯度是 \(\nabla_hv\) 的水平提升。投影在水平空间上保长，故 \(|\mathrm du|_{g_\varepsilon}^2=|\mathrm dv|_h^2\circ\pi\)。体积公式更精确地给 \(\pi_\#\mu_{g_\varepsilon}=2\pi\varepsilon\mu_h\)。除以 \(2\pi\varepsilon\operatorname{Vol}_hB\)，对 \(v^2\) 与 \(|\mathrm dv|_h^2\) 积分即得两式。这个计算表明，纤维方向消失时，沿纤维常值的函数仍保留基底上的质量和能量。
\end{solution}

\begin{exercise}\label{g3:ex:collapse-order-of-scales}
	在 \(T_{1,\varepsilon}\)、\(0<\varepsilon<1\)，上以 \(p\) 为任意基点，记归一化面积为 \(\nu_\varepsilon\)。证明
	\[
		\lim_{\varepsilon\downarrow0}\lim_{r\downarrow0}
		\frac{\nu_\varepsilon(B(p,r))}{2r}=0,
		\qquad
		\lim_{r\downarrow0}\lim_{\varepsilon\downarrow0}
		\frac{\nu_\varepsilon(B(p,r))}{2r}=1.
	\]
\end{exercise}
\begin{solution}
	固定 \(\varepsilon\)，当 \(r<\varepsilon/2\) 时球是 Euclid 圆盘，面积为 \(\pi r^2\)，而总体积为 \(\varepsilon\)。所以内层极限为 \(\lim_{r\downarrow0}\pi r/(2\varepsilon)=0\)。

	反过来，固定 \(0<r<1/2\)，将圆周坐标表示在 \([-1/2,1/2]\) 与 \([-\varepsilon/2,\varepsilon/2]\) 中。当 \(\varepsilon<2r\) 时，乘积距离公式给
	\[
		\nu_\varepsilon(B(p,r))
		=\frac1\varepsilon\int_{-\varepsilon/2}^{\varepsilon/2}
		2\sqrt{r^2-y^2}\,\mathrm dy\longrightarrow2r.
	\]
	因此第二个迭代极限为一。先缩小观察球看到二维切空间，先压缩纤维则看到极限圆周上的一维长度测度。
\end{solution}

\begin{exercise}\label{g3:ex:metric-limit-measure-choice}
	在固定长为 \(2\pi\) 的圆周上，让概率测度在 \(\nu_0=\mathrm dx/(2\pi)\) 与 \(\nu_1=(2+\sin x)\mathrm dx/(4\pi)\) 之间交替。证明它在这个固定空间上没有弱极限，尽管度量始终不变。进一步证明 \((S^1,d,\nu_0)\) 与 \((S^1,d,\nu_1)\) 不存在保持测度的满等距映射。
\end{exercise}
\begin{solution}
	以 \(f(x)=\sin x\) 测试，得到
	\[
		\int f\,\mathrm d\nu_0=0,\qquad
		\int f\,\mathrm d\nu_1=\frac1{4\pi}\int_0^{2\pi}\sin^2x\,\mathrm dx=\frac14.
	\]
	积分交替取两个不同值，故不弱收敛。圆周的满等距映射是旋转或旋转后反射：选定一个点及其一侧的短弧方向后，保距性确定每个点的弧长坐标。两种映射都保持长度测度，因此把 \(\nu_0\) 推前仍得 \(\nu_0\)，不可能得到非恒定密度的 \(\nu_1\)。距离的极限本身因而不能唯一指定附加测度。
\end{solution}

\begin{exercise}\label{g3:ex:full-support-limit}
	在通常区间 \(X=[0,1]\) 上，令 \(\lambda\) 为长度概率测度，并对整数 \(j\ge2\) 定义
	\[
		\nu_j=(1-j^{-1})j\mathbf1_{[0,1/j]}\lambda+j^{-1}\lambda.
	\]
	证明每个 \(\nu_j\) 满支撑，而其弱极限为不满支撑的 \(\delta_0\)。说明\cref{g3:thm:ricci-measured-precompactness}中的几何估计在哪一步排除了这一现象。
\end{exercise}
\begin{solution}
	每个非空相对开集的长度为正，\(\nu_j\ge j^{-1}\lambda\)，故满支撑。对连续函数 \(f\)，一致连续性给
	\[
		j\int_0^{1/j}f(x)\,\mathrm dx\longrightarrow f(0),
	\]
	第二项的质量为 \(j^{-1}\)，因此 \(\int f\,\mathrm d\nu_j\to f(0)\)。极限是 \(\delta_0\)，支撑仅为 \(\{0\}\)。正文的 Ricci 与直径条件提供对任意固定 \(r>0\) 都独立于 \(j\) 和球心的正球质量下界，并利用球测度收敛传到极限。这里远离零的固定小球质量趋零，恰好缺少这个下界。
\end{solution}

\begin{exercise}\label{g3:ex:cone-product-singular-axis}
	设 \(0<\alpha<2\pi\)，\(C_\alpha\) 为角度 \(\alpha\) 的平坦二维度量锥，顶点为 \(o\)。在乘积 \(X=C_\alpha\times\mathbb R^k\) 中，\(k\ge0\)，求各点的切锥及奇异点集，并计算以 \((o,0)\) 为中心的半径 \(r\) 球的 \((k+2)\) 维 Hausdorff 测度。
\end{exercise}
\begin{solution}
	离开 \(\{o\}\times\mathbb R^k\) 的点有局部 Euclid 坐标，所以切锥唯一且为 \(\mathbb R^{k+2}\)。在轴上，先平移 Euclid 分量，再同时伸缩锥半径与 Euclid 坐标，得到保持基点的缩放等距。因此每个轴点的所有切锥都等距于 \(C_\alpha\times\mathbb R^k\)。

	令 \(n=k+2\)。在锥的极坐标中，体积元为 \(\rho\,\mathrm d\rho\,\mathrm d\theta\,\mathrm dz\)，其中 \(0\le\theta<\alpha\)。球条件为 \(\rho^2+|z|^2<r^2\)。与 Euclid 空间相同积分相比，只有角变量范围由 \(2\pi\) 变为 \(\alpha\)，故
	\[
		\mathcal H^n(B((o,0),r))=\frac{\alpha}{2\pi}\omega_nr^n.
	\]
	这里采用与 Euclid 体积一致的 Hausdorff 归一化；锥轴本身的 \(\mathcal H^n\) 测度为零。因 \(\alpha<2\pi\)，此球体积不同于 Euclid 球，故轴点切锥不是 Euclid 空间。奇异集正是 \(\{o\}\times\mathbb R^k\)，其 Hausdorff 维数为 \(k=n-2\)。
\end{solution}

\begin{exercise}\label{g3:ex:two-point-transport}
	在区间 \([0,L]\)、\(L>0\)，上令 \(\mu_a=(1-a)\delta_0+a\delta_L\)，其中 \(0\le a\le1\)。证明 \(W_2(\mu_a,\mu_b)^2=L^2|a-b|\)。比较路径 \(t\mapsto\mu_t\) 与 \(t\mapsto\delta_{tL}\)，判断哪一条是从 \(\delta_0\) 到 \(\delta_L\) 的常速 \(W_2\) 测地线。
\end{exercise}
\begin{solution}
	任何耦合都必须把至少 \(|a-b|\) 的质量从一个端点送到另一个端点，每单位的平方代价为 \(L^2\)。把共同质量留在原地、仅运输这个差额，便达到该下界，因此公式成立。于是 \(W_2(\mu_s,\mu_t)=L\sqrt{|t-s|}\)，不等于所需的 \(L|t-s|\)。两个 Dirac 测度之间只有一个耦合，故 \(W_2(\delta_{sL},\delta_{tL})=L|s-t|\)，第二条路径才是常速测地线。测度的线性混合与质量沿空间道路移动是两种不同的插值。
\end{solution}

\begin{exercise}\label{g3:ex:cheeger-scaling}
	设 \((X,d,\nu)\) 为满支撑的紧度量概率空间，\(a>0\)。用局部斜率松弛的定义证明
	\[
		\operatorname{Ch}_{ad,\nu}(f)=a^{-2}\operatorname{Ch}_{d,\nu}(f)
		\qquad(f\in L^2(\nu)).
	\]
	更一般地，将参考测度同时乘以常数 \(b>0\) 时，公式怎样改变？
\end{exercise}
\begin{solution}
	两个距离给出相同拓扑和相同 Lipschitz 函数类，且 \(\operatorname{lip}_{ad}u=a^{-1}\operatorname{lip}_du\)。参考测度不变，逼近序列的 \(L^2\) 收敛条件相同；因此定义中的每个积分都乘 \(a^{-2}\)，取下极限再取下确界给出第一式。若测度改为 \(b\nu\)，\(L^2\) 范数只乘 \(\sqrt b\)，其收敛序列仍相同，而能量积分乘 \(b\)。故一般公式为 \(\operatorname{Ch}_{ad,b\nu}(f)=ba^{-2}\operatorname{Ch}_{d,\nu}(f)\)，允许 \(b\nu\) 为有限非概率测度。
\end{solution}

\begin{exercise}\label{g3:ex:weighted-circle-ricci}
	将\cref{g3:exm:weighted-circle-collapse}的极限测度写成 \(\nu=Z^{-1}e^{-V}\mathrm dx\)，其中 \(V=-\log(2+\sin x)\)、\(Z=4\pi\)。计算加权 Ricci 张量 \(\operatorname{Ric}+\operatorname{Hess}V\) 的最佳常数下界，并据\cref{g3:thm:smooth-cd-rcd-comparison}判定相应的 \(\mathrm{RCD}(K,\infty)\) 条件。
\end{exercise}
\begin{solution}
	一维圆周的 Ricci 张量为零，Hessian 为 \(V''\mathrm dx^2\)。计算得
	\[
		V'=-\frac{\cos x}{2+\sin x},\qquad
		V''=\frac{1+2\sin x}{(2+\sin x)^2}.
	\]
	令 \(s=\sin x\)，函数 \((1+2s)/(2+s)^2\) 在 \([-1,1]\) 上的导数为 \(2(1-s)/(2+s)^3\ge0\)，故最小值为 \(-1\)，在 \(s=-1\) 处达到。最佳下界为 \(-\mathrm dx^2\)。\Cref{g3:thm:smooth-cd-rcd-comparison}说明，该带权圆周满足 \(\mathrm{RCD}(-1,\infty)\)，也满足所有 \(K\le-1\) 的条件；若 \(K>-1\)，该比较定理的必要方向排除它。这与塌缩前二维曲率的统一下界相容。
\end{solution}

\begin{exercise}\label{g3:ex:quadratic-energy-not-rcd}
	令 \(X=\{0,L\}\)、\(L>0\)，赋予通常距离和两点各质量 \(1/2\) 的概率测度。计算其 Cheeger 能量，证明能量具有二次性；再说明这个空间不满足本章的 \(\mathrm{RCD}(K,\infty)\) 定义。
\end{exercise}
\begin{solution}
	每个点都孤立，按定义局部斜率为零；每个实函数在这个有限空间上又都是 Lipschitz 的。因此以常值逼近序列可得 \(\operatorname{Ch}(f)=0\) 对所有函数成立，平行四边形恒等式成立，Sobolev 空间就是带通常内积的两维 \(L^2\)。

	但 \(X\) 不是测地空间。即使只看概率空间，也没有从 \(\delta_0\) 到 \(\delta_L\) 的常速 \(W_2\) 测地线：任何候选中点都是某个 \(\mu_a\)，而\cref{g3:ex:two-point-transport}给出的两段距离为 \(L\sqrt a\) 与 \(L\sqrt{1-a}\)，不可能同时等于 \(L/2\)。两个端点测度对参考测度都绝对连续且熵有限，所以熵凸性所要求的测地线本身已经不存在。能量二次性只是 RCD 条件的一部分。
\end{solution}

\begin{exercise}\label{g3:ex:c0-metric-eigenvalues}
	设 \((M^n,g_j)\) 在同一个闭光滑流形上以 \(C^0\) 收敛到光滑正定度量 \(g\)。使用非负 Laplace 算子的极小极大表达
	\[
		\lambda_k(g)=\inf_{\substack{E\subset H^1(M)\\\dim E=k+1}}
		\sup_{u\in E\setminus\{0\}}
		\frac{\int_M|\mathrm du|_g^2\mu_g}{\int_Mu^2\mu_g},
		\qquad k=0,1,\ldots,
	\]
	证明每个固定 \(k\) 都有 \(\lambda_k(g_j)\to\lambda_k(g)\)。这里按重数排列，并从零特征值开始编号。
\end{exercise}
\begin{solution}
	正定性与紧性使 \(C^0\) 收敛给出 \(\delta_j\downarrow0\)，满足 \(e^{-2\delta_j}g\le g_j\le e^{2\delta_j}g\)。由\cref{g3:ex:metric-comparison-energy}及体积密度界，任意非零 \(u\in H^1(M)\) 的 Rayleigh 商满足
	\[
		e^{-(2n+2)\delta_j}\mathcal Q_g(u)
		\le\mathcal Q_{g_j}(u)
		\le e^{(2n+2)\delta_j}\mathcal Q_g(u).
	\]
	光滑函数稠密使能量比较延伸到 \(H^1\)；度量的一致等价还说明这些 Sobolev 空间作为函数集合相同。故在同一个有限维子空间取上确界，再对所有 \((k+1)\) 维子空间取下确界，得到相同的两侧特征值界。令 \(j\to\infty\) 即得结论，包括 \(\lambda_k=0\) 的情形。用比较微分同胚拉回后，同样适用于 \(C^0\) Cheeger--Gromov 收敛。
\end{solution}
